\section{Suplemento sobre álgebra homológica}
\label{section:0.11}
\label{section:0.11-fr}

\subsection{Repaso de las sucesiones espectrales}
\label{subsection:0.11.1}
\label{subsection:0.11.1-fr}

\begin{env}[11.1.1]
\label{0.11.1.1-fr}
Utilizaremos en lo que sigue una noción de sucesión espectral más general que
la definida en (T, 2.4); conservando las notaciones de (T, 2.4),
llamaremos \emph{sucesión espectral} en una categoría abeliana $\boldsymbol{C}$ a un
sistema E formado por los elementos siguientes:
\begin{enumerate}
  \item[$a)$] Una familia $(E_r^{pq})$ de objetos de $\boldsymbol{C}$ definidos para
    $p\in\mathbb Z$, $q\in\mathbb Z$ y $r\geq2$.
  \item[$b)$] Una familia de morfismos
    $d_r^{pq}:E_r^{pq}\to E_r^{p+r,q-r+1}$ tales que
    $d_r^{p+r,q-r+1}d_r^{pq}=0$. Se pone
    $Z_{r+1}(E_r^{pq})=\operatorname{Ker}(d_r^{pq})$,
    $B_{r+1}(E_r^{pq})=\operatorname{Im}(d_r^{p+r,q-r+1})$, de modo que
    \[
      B_{r+1}(E_r^{pq})\subset Z_{r+1}(E_r^{pq})\subset E_r^{pq}.
    \]
  \item[$c)$] Una familia de isomorfismos
    $\alpha_r^{pq}:Z_{r+1}(E_r^{pq})/B_{r+1}(E_r^{pq})
    \xrightarrow{\sim}E_{r+1}^{pq}$.
\end{enumerate}

Se definen entonces para $k\geq r+1$, por inducción sobre $k$, los subobjetos
$B_k(E_r^{pq})$ y $Z_k(E_r^{pq})$ como imágenes inversas, por el morfismo
canónico $E_r^{pq}\to E_r^{pq}/B_{r+1}(E_r^{pq})$ de los subobjetos de este
cociente identificados por $\alpha_r^{pq}$ con los subobjetos $B_k(E_{r+1}^{pq})$ y
$Z_k(E_{r+1}^{pq})$ respectivamente. Está claro que se tiene entonces, salvo
isomorfismo,
\[
\label{0.11.1.1.1-fr}
  Z_k(E_r^{pq})/B_k(E_r^{pq})=E_k^{pq}
  \qquad\text{para }k\geq r+1
  \tag{11.1.1.1}
\]
y, si se pone además $B_r(E_r^{pq})=0$ y $Z_r(E_r^{pq})=E_r^{pq}$, se tienen
las relaciones de inclusión
\[
\label{0.11.1.1.2-fr}
  0=B_r(E_r^{pq})\subset B_{r+1}(E_r^{pq})\subset B_{r+2}(E_r^{pq})
  \subset\cdots\subset Z_{r+2}(E_r^{pq})\subset Z_{r+1}(E_r^{pq})
  \subset Z_r(E_r^{pq})=E_r^{pq}.
  \tag{11.1.1.2}
\]

Los otros elementos del dato de E son entonces:
\begin{enumerate}
  \item[$d)$] Dos subobjetos $B_\infty(E_2^{pq})$ y
    $Z_\infty(E_2^{pq})$ de $E_2^{pq}$ tales que se tenga
    $B_\infty(E_2^{pq})\subset Z_\infty(E_2^{pq})$ y, para todo $k\geq2$,
    \[
      B_k(E_2^{pq})\subset B_\infty(E_2^{pq})
      \qquad\text{y}\qquad
      Z_\infty(E_2^{pq})\subset Z_k(E_2^{pq}).
    \]
    Se pone
    \[
    \label{0.11.1.1.3-fr}
      E_\infty^{pq}=Z_\infty(E_2^{pq})/B_\infty(E_2^{pq}).
      \tag{11.1.1.3}
    \]
  \item[$e)$]
    \oldpage[0\textsubscript{III}]{24}
    Una familia $(E^n)$ de objetos de $\boldsymbol{C}$, cada uno provisto de una
    \emph{filtración decreciente} $(F^p(E^n))_{p\in\mathbb Z}$. Se designa
    como de costumbre por $\operatorname{gr}(E^n)$ el objeto graduado asociado al
    objeto filtrado $E^n$, suma directa de los
    $\operatorname{gr}_p(E^n)=F^p(E^n)/F^{p+1}(E^n)$.
  \item[$f)$] Para todo par $(p,q)\in\mathbb Z\times\mathbb Z$, un
    isomorfismo
    $\beta^{pq}:E_\infty^{pq}\xrightarrow{\sim}
    \operatorname{gr}_p(E^{p+q})$.
\end{enumerate}

La familia $(E^n)$, sin las filtraciones, se llama el \emph{término límite}
de la sucesión espectral E.

Supongamos que la categoría $\boldsymbol{C}$ admite sumas directas infinitas,
o que para todo $r\geq2$ y todo $n\in\mathbb Z$ los pares $(p,q)$ tales
que $p+q=n$ y $E_r^{pq}\ne0$ son en número finito (basta que sea así
para $r=2$). Entonces los
\[
  E_r^{(n)}=\sum_{p+q=n}E_r^{pq}
\]
están definidos, y designando por $d_r^{(n)}$ el morfismo
$E_r^{(n)}\to E_r^{(n+1)}$ cuya restricción a $E_r^{pq}$ es $d_r^{pq}$
para cada par $(p,q)$ tal que $p+q=n$, se tiene
$d_r^{(n+1)}\circ d_r^{(n)}=0$, en otros términos $(E_r^{(n)})_{n\in\mathbb Z}$
es un \emph{complejo $E_r^{(\bullet)}$} en $\boldsymbol{C}$, con operador de
derivación de grado $+1$, y resulta de \emph{c)} que
$\mathrm H^n(E_n^{(\bullet)})$ es \emph{isomorfo a} $E_{r+1}^{(n)}$ para todo
$r\geq2$.
\end{env}

\begin{env}[11.1.2]
\label{0.11.1.2}
\label{0.11.1.2-fr}
Un \emph{morfismo $u:E\to E'$} de una sucesión espectral E en una sucesión espectral $E'=(E_r^{\prime pq},E^{\prime n})$ consiste en sistemas de morfismos $u_r^{pq}:E_r^{pq}\to E_r^{\prime pq}$ y $u^n:E^n\to E^{\prime n}$, siendo los $u^n$ compatibles con las filtraciones de $E^n$ y $E^{\prime n}$, y tales que los diagramas
\[
  \xymatrix{
    E_r^{pq}\ar[r]^-{d_r^{pq}}\ar[d]_{u_r^{pq}} &
    E_r^{p+r,q-r+1}\ar[d]^{u_r^{p+r,q-r+1}}\\
    E_r^{\prime pq}\ar[r]_{d_r^{\prime pq}} &
    E_r^{\prime p+r,q-r+1}
  }
\]
sean conmutativos; además, pasando a los cocientes, $u_r^{pq}$ da un morfismo \[
  \overline u_r^{pq}:
  Z_{r+1}(E_r^{pq})/B_{r+1}(E_r^{pq})\longrightarrow
  Z_{r+1}(E_r^{\prime pq})/B_{r+1}(E_r^{\prime pq})
\] y debe tenerse $\alpha_r^{\prime pq}\circ\overline u_r^{pq}
=u_{r+1}^{pq}\circ\alpha_r^{pq}$; finalmente, debe tenerse $u_2^{pq}(B_\infty(E_2^{pq}))\subset B_\infty(E_2^{\prime pq})$ y $u_2^{pq}(Z_\infty(E_2^{pq}))\subset Z_\infty(E_2^{\prime pq})$; pasando a los cocientes, $u_2^{pq}$ da entonces un morfismo $u_\infty^{\prime pq}:E_\infty^{pq}\to E_\infty^{\prime pq}$, y el diagrama
\[
  \xymatrix{
    E_\infty^{pq}\ar[r]^-{u_\infty^{\prime pq}}\ar[d]_{\beta^{pq}} &
    E_\infty^{\prime pq}\ar[d]^{\beta^{\prime pq}}\\
    \operatorname{gr}_{\boldsymbol{p}}(E^{p+q})
      \ar[r]_{\operatorname{gr}_{\boldsymbol{p}}(u^{p+q})} &
    \operatorname{gr}_{\boldsymbol{p}}(E^{\prime\boldsymbol{p}+q})
  }
\]
debe ser conmutativo.

Las definiciones anteriores muestran, por inducción sobre $r$, que si los $u_2^{pq}$ son \emph{isomorfismos}, también lo son los $u_r^{pq}$ para $r\geq2$; \emph{si además se sabe} que $u_2^{pq}(B_\infty(E_2^{pq}))=B_\infty(E_2^{\prime pq})$ y $u_2^{pq}(Z_\infty(E_2^{pq}))=Z_\infty(E_2^{\prime pq})$ y que los $u^n$ son \emph{isomorfismos}, entonces se puede concluir que $u$ es un \emph{isomorfismo}.
\end{env}

\begin{env}[11.1.3]
\label{0.11.1.3}
\label{0.11.1.3-fr}
\oldpage[0\textsubscript{III}]{25}
Recordemos que si $(F^p(X))_{p\in\mathbb Z}$ es una \emph{filtración} (decreciente) de un objeto $X\in\boldsymbol{C}$, se dice que esta filtración es \emph{separada} si $\inf(F^p(X))=0$, \emph{discreta} si existe un $p$ tal que $F^p(X)=0$, \emph{exhaustiva} (o \emph{coseparada}) si $\sup(F^p(X))=X$, y \emph{codiscreta} si existe un $p$ tal que $F^p(X)=X$.

Diremos que una sucesión espectral $E=(E_r^{pq},E^n)$ es \emph{débilmente convergente} si se tiene $B_\infty(E_2^{pq})=\sup_k(B_k(E_2^{pq}))$ y $Z_\infty(E_2^{pq})=\inf_k(Z_k(E_2^{pq}))$ (en otras palabras, los objetos $B_\infty(E_2^{pq})$ y $Z_\infty(E_2^{pq})$ están determinados por los datos \emph{a)} a \emph{c)} de la sucesión espectral E). Diremos que la sucesión espectral E es \emph{regular} si es débilmente convergente y si además:
\begin{enumerate}
  \item[$1^\circ$] Para todo par $(p,q)$, la sucesión decreciente $(Z_k(E_2^{pq}))_{k\geq2}$ es \emph{estable}; la hipótesis de que E es débilmente convergente implica entonces
$Z_\infty(E_2^{pq})=Z_k(E_2^{pq})$ para $k$ suficientemente grande (dependiendo de $p$ y $q$).
  \item[$2^\circ$] Para todo $n$, la filtración $(F^p(E^n))_{p\in\mathbb Z}$ de $E^n$ es \emph{discreta} y \emph{exhaustiva}.
\end{enumerate}

Se dice que la sucesión espectral E es \emph{corregular} si es débilmente convergente y si además:
\begin{enumerate}
  \item[$3^\circ$] Para todo par $(p,q)$, la sucesión creciente $(B_k(E_2^{pq}))_{k\geq2}$ es \emph{estable}, lo que implica que $B_\infty(E_2^{pq})=B_k(E_2^{pq})$ y, en consecuencia, $E_\infty^{pq}=\inf E_k^{pq}$.
  \item[$4^\circ$] Para todo $n$, la filtración de $E^n$ es \emph{codiscreta}.
\end{enumerate}

Finalmente, se dice que E es \emph{birregular} si es a la vez regular y corregular, es decir, si se tienen las condiciones siguientes:
\begin{enumerate}
  \item[\textit{a)}] Para todo par $(p,q)$, las sucesiones $(B_k(E_2^{pq}))_{k\geq2}$ y $(Z_k(E_2^{pq}))_{k\geq2}$ son \emph{estables} y se tiene $B_\infty(E_2^{pq})=B_k(E_2^{pq})$ y $Z_\infty(E_2^{pq})=Z_k(E_2^{pq})$ para $k$ suficientemente grande (lo que implica que $E_\infty^{pq}=E_k^{pq}$).
  \item[\textit{(b)}] Para todo $n$, la filtración $(F^p(E^n))_{p\in\mathbb Z}$ es \emph{discreta} y \emph{codiscreta} (también se dice \emph{finita}).
\end{enumerate}

Las sucesiones espectrales definidas en (T, 2.4) son, por tanto, sucesiones espectrales birregulares.
\end{env}

\begin{env}[11.1.4]
\label{0.11.1.4}
\label{0.11.1.4-fr}
Supongamos que en la categoría $\boldsymbol{C}$ existen los límites inductivos filtrantes y que el funtor $\varinjlim$ es \emph{exacto} (lo que equivale a decir que se satisface el axioma (AB~5) de (T, 1.5) (cf.~T, 1.8)).
La condición de que la filtración $(F^p(X))_{p\in\mathbb Z}$ de un objeto $X\in\boldsymbol{C}$ sea exhaustiva se expresa entonces por $\varinjlim_{p\to-\infty}F^p(X)=X$. Si una sucesión espectral E es débilmente convergente, se tiene
$B_\infty(E_2^{pq})=\varinjlim_{k\to\infty}B_k(E_2^{pq})$; si además $u:E\to E'$ es un morfismo de E en una sucesión espectral débilmente convergente $E'$ en $\boldsymbol{C}$, entonces se tiene $u_2^{pq}(B_\infty(E_2^{pq}))=B_\infty(E_2^{\prime pq})$, por la exactitud de $\varinjlim$.
Además:
\end{env}

\begin{proposition}[11.1.5]
\label{0.11.1.5}
\label{0.11.1.5-fr}
Sea $\mathbf{C}$ una categoría abeliana en la que los límites
inductivos filtrantes son exactos, \textup{E}, $E'$ dos sucesiones espectrales regulares en $\mathbf{C}$, y sea $u:E\to E'$ un morfismo de sucesiones espectrales.
Si los $u_2^{pq}$ son isomorfismos, entonces $u$ también lo es.
\end{proposition}

\begin{proof}
Ya sabemos por \sref{0.11.1.2} que los $u_r^{pq}$ son isomorfismos y que
\[
  u_2^{pq}(B_\infty(E_2^{pq}))=B_\infty(E_2^{\prime pq})~;
\]
\oldpage[0\textsubscript{III}]{26}
la hipótesis de que E y $E'$ son regulares implica también que $u_2^{pq}(Z_\infty(E_2^{pq}))=Z_\infty(E_2^{\prime pq})$ y, como $u_2^{pq}$ es un isomorfismo, también lo es $u_\infty^{pq}$; se concluye así que $\operatorname{gr}_p(u^{p+q})$ es asimismo un isomorfismo.
Pero, como las filtraciones de los $E^n$ y los $E^{\prime n}$ son discretas y exhaustivas, esto implica que los $u^n$ son también isomorfismos (Bourbaki, \emph{Alg. comm.}, cap.~III, \textsection2, n\textsuperscript{o}8, th.~1).
\end{proof}

\begin{env}[11.1.6]
\label{0.11.1.6-fr}
Resulta de (\hyperref[0.11.1.1.2-fr]{11.1.1.2}) y de la definición (\hyperref[0.11.1.1.3-fr]{11.1.1.3}) que si, en una sucesión espectral E, se tiene
$E_r^{pq}=0$, entonces $E_k^{pq}=0$ para $k\geq r$ y $E_\infty^{pq}=0$.
Se dice que una sucesión espectral \emph{degenera} si existe un entero $r\geq2$ y, para todo entero $n\in\mathbb Z$, un entero $q(n)$ tal que $E_r^{n-q,q}=0$ \emph{para todo} $q\ne q(n)$.
De la observación anterior se deduce primero que también se tiene $E_k^{n-q,q}=0$ para $k\geq r$ (incluido $k=\infty$) y $q\ne q(n)$.
Además, la definición de $E_{r+1}^{pq}$ muestra que se tiene $E_{r+1}^{n-q(n),q(n)}=E_r^{n-q(n),q(n)}$; si E es \emph{débilmente convergente}, se tiene, pues, también
$E_\infty^{n-q(n),q(n)}=E_r^{n-q(n),q(n)}$; en otras palabras, para todo $n\in\mathbb Z$, $\operatorname{gr}_p(E^n)=0$ para $p\ne q(n)$ y $\operatorname{gr}_{q(n)}(E^n)=E_r^{n-q(n),q(n)}$.
Si además la filtración de $E^n$ es \emph{discreta} y \emph{exhaustiva}, la sucesión E es \emph{regular} y se tiene
$E^n=E_r^{n-q(n),q(n)}$ salvo isomorfismo.
\end{env}

\begin{env}[11.1.7]
\label{0.11.1.7}
\label{0.11.1.7-fr}
Supongamos que en la categoría $\boldsymbol{C}$ existen los límites inductivos filtrantes y son exactos, y sea $(E_\lambda,u_{\mu\lambda})$ un sistema inductivo (sobre un conjunto filtrante de índices) de sucesiones espectrales en $\boldsymbol{C}$.
Entonces el \emph{límite inductivo} de este sistema inductivo existe en la categoría aditiva de las sucesiones espectrales de objetos de $\boldsymbol{C}$: para verlo, basta definir $E_r^{pq}$, $d_r^{pq}$, $\alpha_r^{pq}$, $B_\infty(E_2^{pq})$, $Z_\infty(E_2^{pq})$, $E^n$, $F^p(E^n)$ y $\beta^{pq}$ como los límites inductivos respectivos de los $E_{r,\lambda}^{pq}$, $d_{r,\lambda}^{pq}$, $\alpha_{r,\lambda}^{pq}$, $B_\infty(E_{2,\lambda}^{pq})$, $Z_\infty(E_{2,\lambda}^{pq})$, $E_\lambda^n$, $F^p(E_\lambda^n)$ y $\beta_\lambda^{pq}$; la verificación de las condiciones de \sref{0.11.1.1} se deduce de la exactitud del funtor $\varinjlim$ sobre $\boldsymbol{C}$.
\end{env}

\begin{remark}[11.1.8]
\label{0.11.1.8}
\label{0.11.1.8-fr}
Supongamos que la categoría $\boldsymbol{C}$ es la categoría de los $A$-módulos sobre un anillo $A$ \emph{noetheriano} (resp. un anillo $A$).
Las definiciones de \sref{0.11.1.1} muestran entonces que si, para un $r$ dado, los $E_r^{pq}$ son $A$-módulos \emph{de tipo finito} (resp. \emph{de longitud finita}), también lo son todos los módulos $E_s^{pq}$ para $s\geq r$ y, por tanto, también $E_\infty^{pq}$.
Si además la filtración del término límite $(E^n)$ es \emph{discreta} y \emph{codiscreta} para todo $n$, se concluye que cada $E_n$ es también un $A$-módulo \emph{de tipo finito} (resp. \emph{de longitud finita}).
\end{remark}

\begin{env}[11.1.9]
\label{0.11.1.9-fr}
Habremos de considerar condiciones que aseguren que una sucesión espectral E sea birregular de manera «uniforme» en $p+q=n$.
Se utilizará entonces el lema siguiente:
\end{env}

\begin{lemma}[11.1.10]
\label{0.11.1.10}
\label{0.11.1.10-fr}
Sea $(E_r^{pq})$ una familia de objetos de $\mathbf{C}$ relacionados por los datos \emph{(a)}, \emph{(b)} y \emph{(c)} de \textup{\sref{0.11.1.1}}.
Para un entero fijo $n$, las propiedades siguientes son equivalentes:
\begin{enumerate}
  \item[{\rm(a)}] Existe un entero $r(n)$ tal que, para $r\geq r(n)$, $p+q=n$ o $p+q=n-1$, los morfismos $d_r^{pq}$ son nulos.
  \item[{\rm(b)}] Existe un entero $r(n)$ tal que, para $p+q=n$ o $p+q=n+1$, se tiene $B_r(E_2^{pq})=B_s(E_2^{pq})$ para $s\geq r\geq r(n)$.
  \item[{\rm(c)}] Existe un entero $r(n)$ tal que, para $p+q=n$ o $p+q=n-1$, se tiene $Z_r(E_2^{pq})=Z_s(E_2^{pq})$ para $s\geq r\geq r(n)$.
  \item[{\rm(d)}] Existe un entero $r(n)$ tal que, para $p+q=n$, se tiene $B_r(E_2^{pq})=B_s(E_2^{pq})$ y $Z_r(E_2^{pq})=Z_s(E_2^{pq})$ para $s\geq r\geq r(n)$.
\end{enumerate}
\end{lemma}

\begin{proof}
\oldpage[0\textsubscript{III}]{27}
Según las condiciones \emph{(a), (b) y (c)} de \sref{0.11.1.1}, decir que $Z_{r+1}(E_2^{pq})=Z_r(E_2^{pq})$ equivale a decir que $d_r^{pq}=0$, y decir que $B_r(E_2^{p+r,q-r+1})=B_{r+1}(E_2^{p+r,q-r+1})$ equivale a decir que $d_r^{pq}=0$; el lema se deduce inmediatamente de esta observación.
\end{proof}

\subsection{La sucesión espectral de un complejo filtrado}
\label{subsection:0.11.2}
\label{subsection:0.11.2-fr}

\begin{env}[11.2.1]
\label{0.11.2.1}
\label{0.11.2.1-fr}
Sea $\boldsymbol{C}$ una categoría abeliana; convenimos en denotar mediante notaciones como $K^\bullet$ los \emph{complejos $(K^i)_{i\in\mathbb Z}$} de objetos de $\boldsymbol{C}$ cuyo diferencial es de grado $+1$, y mediante notaciones como $K_\bullet$ los complejos $(K_i)_{i\in\mathbb Z}$ de objetos de $\boldsymbol{C}$ cuyo diferencial es de grado $-1$.
A todo complejo $K^\bullet=(K^i)$ cuyo diferencial $d$ es de grado $+1$ se puede asociarle un complejo $K'_\bullet=(K'_i)$ poniendo $K'_i=K^{-i}$, siendo el diferencial $K'_i\to K'_{i-1}$ el operador $d:K^{-i}\to K^{-i+1}$; y \emph{viceversa}, lo que permitirá, según las circunstancias, considerar uno u otro de los tipos de complejos y traducir todo resultado de uno de los tipos en resultados para el otro.
Denotaremos análogamente mediante notaciones como $K^{\bullet\bullet}=(K^{ij})$ (resp. $K_{\bullet\bullet}=(K_{ij})$) los \emph{bicomplejos} (o \emph{complejos dobles}) de objetos de $\boldsymbol{C}$ en los que los \emph{dos} diferenciales son de grado $+1$ (resp. $-1$); todavía se puede pasar de un tipo al otro cambiando los signos de los índices, y se tienen notaciones y observaciones análogas para los multicomplejos cualesquiera.
Las notaciones $K^\bullet$ y $K_\bullet$ se utilizará también para \emph{objetos graduados} de $\boldsymbol{C}$, de tipo $\mathbb Z$, que no son necesariamente complejos (o que se pueden considerar como tales para operadores de derivación \emph{nulos}); por ejemplo, escribiremos
$\mathrm H^\bullet(K^\bullet)=(\mathrm H^i(K^\bullet))_{i\in\mathbb Z}$ para la \emph{cohomología} de un complejo $K^\bullet$ cuyo diferencial es de grado $+1$, y $\mathrm H_\bullet(K_\bullet)=(\mathrm H_i(K_\bullet))_{i\in\mathbb Z}$ para la \emph{homología} de un complejo $K_\bullet$ cuyo diferencial es de grado $-1$; cuando se pasa de $K^\bullet$ a $K'_\bullet$ por el procedimiento descrito anteriormente, se tiene $\mathrm H_i(K'_\bullet)=\mathrm H^{-i}(K^\bullet)$.

Recordemos en este caso que, para un complejo $K^\bullet$ (resp. $K_\bullet$), escribiremos en general $Z^i(K^\bullet)=\operatorname{Ker}(K^i\to K^{i+1})$ (``objeto de cociclos'') y $B^i(K^\bullet)=\operatorname{Im}(K^{i-1}\to K^i)$ (``objeto de cobordes'') (resp. $Z_i(K_\bullet)=\operatorname{Ker}(K_i\to K_{i-1})$ (``objeto de ciclos'') y $B_i(K_\bullet)=\operatorname{Im}(K_{i+1}\to K_i))$ (``objeto de bordes'')), de modo que $\mathrm H^i(K^\bullet)=Z^i(K^\bullet)/B^i(K^\bullet)$ (resp. $\mathrm H_i(K_\bullet)=Z_i(K_\bullet)/B_i(K_\bullet)$).

Si $K^\bullet=(K^i)$ (resp. $K_\bullet=(K_i)$) es un complejo en $\boldsymbol{C}$ y $T:\boldsymbol{C}\to\boldsymbol{C}'$ un funtor de $\boldsymbol{C}$ en una categoría abeliana $\boldsymbol{C}'$, denotaremos por $T(K^\bullet)$ (resp. $T(K_\bullet)$) el complejo $(T(K^i))$ (resp. $(T(K_i))$) en $\boldsymbol{C}'$.

No repasaremos la definición de los \emph{$\partial$-funtores} (T, 2.1), salvo para observar que diremos \emph{también} $\partial$-funtor en lugar de $\partial^*$-funtor cuando el morfismo $\partial$ disminuye el grado en una unidad, aclarando el contexto este punto si puede haber confusión.

Finalmente, diremos que un \emph{objeto graduado $(A_i)_{i\in\mathbb Z}$} de $\boldsymbol{C}$ está \emph{acotado inferiormente} (resp. \emph{superiormente}) si existe un $i_0$ tal que $A_i=0$ para $i<i_0$ (resp. $i>i_0$).
\end{env}

\begin{env}[11.2.2]
\label{0.11.2.2}
\label{0.11.2.2-fr}
Sea $K^\bullet$ un complejo en $\boldsymbol{C}$ cuyo diferencial $d$ es de grado $+1$, y supongamos que está provisto de una \emph{filtración $F(K^\bullet)=(F^p(K^\bullet))_{p\in\mathbb Z}$} formada por subobjetos \emph{graduados}
\oldpage[0\textsubscript{III}]{28}
de $K^\bullet$, es decir, $F^p(K^\bullet)=(K^i\cap F^p(K^\bullet))_{i\in\mathbb Z}$; se supone además que $d(F^p(K^\bullet))\subset F^p(K^\bullet)$ para todo $p\in\mathbb Z$.
Recordemos rápidamente cómo se define \emph{funtorialmente} una sucesión espectral $E(K^\bullet)$ a partir de $K^\bullet$ (M,~XV,~4 y G,~I,~4.3).
Para $r\geq2$, el morfismo canónico $F^p(K^\bullet)/F^{p+r}(K^\bullet)\to
F^p(K^\bullet)/F^{p+1}(K^\bullet)$ define un morfismo en cohomología
\[
  \mathrm H^{p+q}(F^p(K^\bullet)/F^{p+r}(K^\bullet))
  \longrightarrow
  \mathrm H^{p+q}(F^p(K^\bullet)/F^{p+1}(K^\bullet))
\]

Se denota por $Z_r^{pq}(K^\bullet)$ la imagen de este morfismo.
Análogamente, de la sucesión exacta
\[
  0\to F^p(K^\bullet)/F^{p+1}(K^\bullet)
  \to F^{p-r+1}(K^\bullet)/F^{p+1}(K^\bullet)
  \to F^{p-r+1}(K^\bullet)/F^p(K^\bullet)\to0
\]
se deduce de la sucesión exacta de cohomología un morfismo
\[
  \mathrm H^{p+q-1}(F^{p-r+1}(K^\bullet)/F^p(K^\bullet))
  \longrightarrow
  \mathrm H^{p+q}(F^p(K^\bullet)/F^{p+1}(K^\bullet))
\]
y se denota por $B_r^{pq}(K^\bullet)$ la imagen de este morfismo; se demuestra que $B_r^{pq}(K^\bullet)\subset Z_r^{pq}(K^\bullet)$ y se toma $E_r^{pq}(K^\bullet)=Z_r^{pq}(K^\bullet)/B_r^{pq}(K^\bullet)$; no precisaremos la definición de los $d_r^{pq}$ ni de los $\alpha_r^{pq}$. Se observará aquí que todos los $Z_r^{pq}(K^\bullet)$ y $B_p^{pq}(K^\bullet)$, para $p$ y $q$ fijos, son subobjetos del mismo objeto $\mathrm H^{p+q}(F^p(K^\bullet)/F^{p+1}(K^\bullet))$, que se denota por $Z_1^{pq}(K^\bullet)$; se pone $B_1^{pq}(K^\bullet)=0$, de modo que las definiciones anteriores de $Z_r^{pq}(K^\bullet)$ y $B_r^{pq}(K^\bullet)$ valen también para $r=1$; se pone asimismo $E_1^{pq}(K^\bullet)=Z_1^{pq}(K^\bullet)$. Se definen también los $d_1^{pq}$ y $\alpha_1^{pq}$ de modo que se satisfagan las condiciones de \sref{0.11.1.1} para $r=1$.
Por otra parte, se definen los subobjetos $Z_\infty^{pq}(K^\bullet)$ como la imagen del morfismo
\[
  \mathrm H^{p+q}(F^p(K^\bullet))\longrightarrow
  \mathrm H^{p+q}(F^p(K^\bullet)/F^{p+1}(K^\bullet))
  =E_1^{pq}(K^\bullet)
\]
y $B_\infty^{pq}(K^\bullet)$ como la imagen del morfismo
\[
  \mathrm H^{p+q-1}(K^\bullet/F^p(K^\bullet))\longrightarrow
  \mathrm H^{p+q}(F^p(K^\bullet)/F^{p+1}(K^\bullet))
  =E_1^{pq}(K^\bullet)
\] deducido como arriba de una sucesión exacta de cohomología. Se toma como $Z_\infty(E_2^{pq}(K^\bullet))$ y $B_\infty(E_2^{pq}(K^\bullet))$ como las imágenes canónicas en $E_2^{pq}(K^\bullet)$ de $Z_\infty^{pq}(K^\bullet)$ y $B_\infty^{pq}(K^\bullet)$.

Finalmente, se denota por $F^p(\mathrm H^n(K^\bullet))$ la imagen en $\mathrm H^n(K^\bullet)$ del morfismo $\mathrm H^n(F^p(K^\bullet))\to\mathrm H^n(K^\bullet)$ inducido por la inyección canónica $F^p(K^\bullet)\to K^\bullet$; por la sucesión exacta de cohomología, este es también el núcleo del morfismo $\mathrm H^n(K^\bullet)\to\mathrm H^n(K^\bullet/F^p(K^\bullet))$.
Esto define una filtración sobre $E^n(K^\bullet)=\mathrm H^n(K^\bullet)$; no daremos aquí la definición de los isomorfismos $\beta^{pq}$.
\end{env}

\begin{env}[11.2.3]
\label{0.11.2.3}
\label{0.11.2.3-fr}
El carácter \emph{funtorial} de $E(K^\bullet)$ se entiende del modo siguiente: dados dos complejos \emph{filtrados} $K^\bullet$ y $K^{\prime\bullet}$ en $\boldsymbol{C}$ y un morfismo de complejos $u:K^\bullet\to K^{\prime\bullet}$ \emph{compatible con las filtraciones}, se inducen de manera evidente los morfismos $u_r^{pq}$ (para $r\geq1$) y $u^n$, y se demuestra que estos morfismos son compatibles con los $d_r^{pq}$, $\alpha_r^{pq}$ y $\beta^{pq}$ en el sentido de \sref{0.11.1.2}, y definen así un morfismo bien definido $E(u):E(K^\bullet)\to E(K^{\prime\bullet})$ de sucesiones espectrales.
Además, se demuestra que si $u$ y $v$ son morfismos $K^\bullet\to K^{\prime\bullet}$ del tipo anterior, \emph{homótopos de orden $\leq k$}, entonces $u_r^{pq}=v_r^{pq}$ para $r>k$ y $u^n=v^n$ para todo $n$ (M,~XV,~3.1).
\end{env}

\begin{env}[11.2.4]
\label{0.11.2.4-fr}
\oldpage[0\textsubscript{III}]{29}
Supongamos que en $\boldsymbol{C}$ los límites inductivos filtrantes sean exactos. Entonces, si la filtración $(F^p(K^\bullet))$ de $K^\bullet$ es \emph{exhaustiva}, también lo es la filtración $(F^p(\mathrm H^n(K^\bullet)))$ para todo $n$, pues por hipótesis se tiene $K^\bullet=\varinjlim_{p\to-\infty}F^p(K^\bullet)$ y la hipótesis sobre $\boldsymbol{C}$ implica que la cohomología conmuta con los límites inductivos.
Además, por la misma razón, se tiene $B_\infty(E_2^{pq}(K^\bullet))=\sup_k B_k(E_2^{pq}(K^\bullet))$.
Se dice que la filtración $(F^p(K^\bullet))$ de $K^\bullet$ es \emph{regular} si para todo $n$ existe un entero $u(n)$ tal que $\mathrm H^n(F^p(K^\bullet))=0$ para $p>u(n)$.
Este es, en particular, el caso cuando la filtración de $K^\bullet$ es \emph{discreta}.
Cuando la filtración de $K^\bullet$ es regular y exhaustiva, y los límites inductivos filtrantes son exactos en $\boldsymbol{C}$, se tiene (M,~XV,~4) que la sucesión espectral $E(K^\bullet)$ es \emph{regular}.
\end{env}

\subsection{Las sucesiones espectrales de un bicomplejo}
\label{subsection:0.11.3}
\label{subsection:0.11.3-fr}

\begin{env}[11.3.1]
\label{0.11.3.1}
\label{0.11.3.1-fr}
En cuanto a las convenciones para los bicomplejos, seguimos las de~(T,~2.4) antes que las de~(M), suponiéndose así que los dos diferenciales $d'$ y $d''$ (de grado~$+1$) de un tal bicomplejo $K^{\bullet\bullet}=(K^{ij})$ se suponen, pues, \emph{permutables}. Supongamos satisfecha \emph{una} de las dos condiciones siguientes: $1^\circ$ Las \emph{sumas directas infinitas} existen en $\boldsymbol{C}$; $2^\circ$ Para todo $n\in\mathbb Z$, solo hay un número \emph{finito} de pares $(p,q)$ tales que $p+q=n$ y $K^{pq}\ne0$.
 
Entonces el bicomplejo $K^{\bullet\bullet}$ define un \emph{complejo} (simple) $(K^{\prime n})_{n\in\mathbb Z}$ con $K^{\prime n}=\sum_{i+j=n}K^{ij}$, viniendo dado el diferencial $d$ (de grado~$+1$) de este complejo por $dx=d'x+(-1)^i d''x$ para $x\in K^{ij}$.
Cuando hablemos más adelante del \emph{complejo (simple) definido por un bicomplejo $K^{\bullet\bullet}$}, se sobreentenderá siempre que se satisface una de las condiciones anteriores.
Se adoptarán convenciones análogas para los multicomplejos.

Se denota por $K^{i,\bullet}$ (resp.~$K^{\bullet,j}$) el complejo simple $(K^{ij})_{j\in\mathbb Z}$ (resp.~$(K^{ij})_{i\in\mathbb Z}$), y por $Z_{\mathrm{II}}^p(K^{i,\bullet})$, $B_{\mathrm{II}}^p(K^{i,\bullet})$, $\mathrm H_{\mathrm{II}}^p(K^{i,\bullet})$ (resp.~$Z_{\mathrm I}^p(K^{\bullet,j})$, $B_{\mathrm I}^p(K^{\bullet,j})$, $\mathrm H_{\mathrm I}^p(K^{\bullet,j})$) sus objetos de cociclos, de cobordes y de cohomología de grado $p$, respectivamente;
el diferencial $d':K^{i,\bullet}\to K^{i+1,\bullet}$ es un morfismo de complejos, que da así un operador sobre los cociclos, los cobordes y la cohomología,
\[
\begin{aligned}
  d'&:Z_{\mathrm{II}}^p(K^{i,\bullet})
      \longrightarrow Z_{\mathrm{II}}^p(K^{i+1,\bullet})\\
  d'&:B_{\mathrm{II}}^p(K^{i,\bullet})
      \longrightarrow B_{\mathrm{II}}^p(K^{i+1,\bullet})\\
  d'&:\mathrm H_{\mathrm{II}}^p(K^{i,\bullet})
      \longrightarrow \mathrm H_{\mathrm{II}}^p(K^{i+1,\bullet})
\end{aligned}
\]
y está claro que, para estos operadores, $(Z_{\mathrm{II}}^p(K^{i,\bullet}))_{i\in\mathbb Z}$, $(B_{\mathrm{II}}^p(K^{i,\bullet}))_{i\in\mathbb Z}$ y $(\mathrm H_{\mathrm{II}}^p(K^{i,\bullet}))_{i\in\mathbb Z}$ son complejos;
denotaremos el complejo $(\mathrm H_{\mathrm{II}}^p(K^{i,\bullet}))_{i\in\mathbb Z}$ por $\mathrm H_{\mathrm{II}}^p(K^{\bullet\bullet})$, y sus objetos de cociclos, cobordes y cohomología de grado $q$ por $Z_{\mathrm I}^q(\mathrm H_{\mathrm{II}}^p(K^{\bullet\bullet}))$, $B_{\mathrm I}^q(\mathrm H_{\mathrm{II}}^p(K^{\bullet\bullet}))$ y $\mathrm H_{\mathrm I}^q(\mathrm H_{\mathrm{II}}^p(K^{\bullet\bullet}))$.
Se definen análogamente los complejos $\mathrm H_{\mathrm I}^p(K^{\bullet\bullet})$ y sus objetos de cohomología $\mathrm H_{\mathrm{II}}^q(\mathrm H_{\mathrm I}^p(K^{\bullet\bullet}))$.
Recordemos, no obstante, que $\mathrm H^n(K^{\bullet\bullet})$ denota el objeto de grado $n$ de la cohomología del complejo \emph{(simple)} definido por $K^{\bullet\bullet}$.
\end{env}

\begin{env}[11.3.2]
\label{0.11.3.2}
\label{0.11.3.2-fr}
Sobre el complejo definido por un bicomplejo $K^{\bullet\bullet}$ pueden considerarse dos filtraciones canónicas, $(F_{\mathrm I}^p(K^{\bullet\bullet}))$ y $(F_{\mathrm{II}}^p(K^{\bullet\bullet}))$, dadas por
\[
\label{0.11.3.2.1-fr}
  F_{\mathrm I}^p(K^{\bullet\bullet})
  =\left(\sum_{i+j=n,\,i\geq p}K^{ij}\right)_{n\in\mathbb Z},
  \qquad
  F_{\mathrm{II}}^p(K^{\bullet\bullet})
  =\left(\sum_{i+j=n,\,j\geq p}K^{ij}\right)_{n\in\mathbb Z}
  \tag{11.3.2.1}
\]
\oldpage[0\textsubscript{III}]{30}
que, por definición, son subobjetos graduados del complejo (simple) definido por $K^{\bullet\bullet}$ y hacen así de este complejo un complejo filtrado;
además, está claro que estas filtraciones son \emph{exhaustivas} y \emph{separadas}.

A cada una de estas filtraciones le corresponde una sucesión espectral~\sref{0.11.2.2};
denotaremos por $'E(K^{\bullet\bullet})$ y $''E(K^{\bullet\bullet})$ las sucesiones espectrales correspondientes a $(F_{\mathrm I}^p(K^{\bullet\bullet}))$ y $(F_{\mathrm{II}}^p(K^{\bullet\bullet}))$, respectivamente, llamadas \emph{sucesiones espectrales del bicomplejo $K^{\bullet\bullet}$}, y que tienen ambas por término límite la cohomología $(\mathrm H^n(K^{\bullet\bullet}))$.
Se demuestra además~(M,~XV,~6) que se tiene
\[
\label{0.11.3.2.2-fr}
  'E_2^{pq}(K^{\bullet\bullet})
  =\mathrm H_{\mathrm I}^p(\mathrm H_{\mathrm{II}}^p(K^{\bullet\bullet})),
  \qquad
  ''E_2^{pq}(K^{\bullet\bullet})
  =\mathrm H_{\mathrm{II}}^p(\mathrm H_{\mathrm I}^p(K^{\bullet\bullet})).
  \tag{11.3.2.2}
\]

Todo morfismo $u:K^{\bullet\bullet}\to K^{\prime\bullet\bullet}$ de bicomplejos es \emph{ipso facto} compatible con las filtraciones del mismo tipo de $K^{\bullet\bullet}$ y $K^{\prime\bullet\bullet}$, y define así un morfismo para cada una de las dos sucesiones espectrales;
además, dos morfismos \emph{homótopos} definen una homotopía \emph{de orden $\leq1$} de los complejos (simples) filtrados correspondientes, y por tanto el \emph{mismo} morfismo para cada una de las dos sucesiones espectrales~(M,~XV,~6.1).
\end{env}

\begin{proposition}[11.3.3]
\label{0.11.3.3}
\label{0.11.3.3-fr}
Sea $K^{\bullet\bullet}=(K^{ij})$ un bicomplejo en una categoría abeliana $\boldsymbol{C}$.
\begin{enumerate}
  \item[{\rm(i)}] Si existen $i_0$ y $j_0$ tales que $K^{ij}=0$ para $i<i_0$ o $j<j_0$ \textup{(resp.}~$i>i_0$ o $j>j_0$), entonces las dos sucesiones espectrales $'E(K^{\bullet\bullet})$ y $''E(K^{\bullet\bullet})$ son birregulares.
  \item[{\rm(ii)}] Si existen $i_0$ e $i_1$ tales que $K^{ij}=0$ para $i<i_0$ o $i>i_1$ \textup{(resp.}~si existen $j_0$ y $j_1$ tales que $K^{ij}=0$ para $j<j_0$ o $j>j_1$), entonces las dos sucesiones espectrales $'E(K^{\bullet\bullet})$ y $''E(K^{\bullet\bullet})$ son birregulares.
\end{enumerate}
Supongamos además que en $\boldsymbol{C}$ los límites inductivos filtrantes existen y son exactos. Entonces:
\begin{enumerate}
\item[\textup{(iii)}] Si existe $i_0$ tal que $K^{ij}=0$ para $i>i_0$ \textup{(resp.}~si existe $j_0$ tal que $K^{ij}=0$ para $j<j_0$), entonces la sucesión espectral $'E(K^{\bullet\bullet})$ es regular.
  \item[{\rm(iv)}] Si existe $i_0$ tal que $K^{ij}=0$ para $i<i_0$ \textup{(resp.}~si existe $j_0$ tal que $K^{ij}=0$ para $j>j_0$), entonces la sucesión espectral $''E(K^{\bullet\bullet})$ es regular.
\end{enumerate}
\end{proposition}

\begin{proof}
La proposición se deduce inmediatamente de las definiciones~\sref{0.11.1.3} y de~\sref{0.11.2.4}, así como de las observaciones siguientes relativas a la filtración $F_{\mathrm I}$ (y de las observaciones análogas que se deducen para $F_{\mathrm{II}}$ intercambiando los papeles de los dos índices en $K^{\bullet\bullet}$):
\begin{enumerate}
  \item[{$1^\circ$}] Si existe $i_0$ tal que $K^{ij}=0$ para $i>i_0$, entonces la filtración $F_{\mathrm I}(K^{\bullet\bullet})$ es \emph{discreta}.
  \item[{$2^\circ$}] Si existe $i_0$ tal que $K^{ij}=0$ para $i<i_0$, entonces la filtración $F_{\mathrm I}(K^{\bullet\bullet})$ es \emph{codiscreta}.
    Se deduce inmediatamente que ocurre lo mismo para la filtración correspondiente $F_{\mathrm I}(\mathrm H^n(K^{\bullet\bullet}))$ para todo $n$.
    Además, la definición de los $B_r^{pq}$ correspondientes a la filtración $F_{\mathrm I}(K^{\bullet\bullet})$ \sref{0.11.2.2} muestra que, para todo par $(p,q)$, la sucesión $(B_r^{pq})_{r\geq2}$ es estable.
  \item[{$3^\circ$}] Si existe $j_0$ tal que $K^{ij}=0$ para $j<j_0$, entonces se tiene
    \[
      F_{\mathrm I}^{p+r}(K^{\bullet\bullet})
      \cap\left(\sum_{i+j=n}K^{ij}\right)=0
    \]
    desde que $p+r+j_0>n$, luego $Z_r^{pq}=Z_\infty(E_2^{pq})$ para $r>q-j_0+1$.
    Por otra parte, $\mathrm H^n(F_{\mathrm I}^p(K^{\bullet\bullet}))=0$ para $p>n-j_0+1$.
  \item[{$4^\circ$}] Si existe $j_0$ tal que $K^{ij}=0$ para $j>j_0$, entonces se tiene
    \[
      F_{\mathrm I}^{p-r+1}(K^{\bullet\bullet})
      \cap\left(\sum_{i+j=n}K^{ij}\right)=\sum_{i+j=n}K^{ij}
    \]
    \oldpage[0\textsubscript{III}]{31}
    desde que $p-r+1+j_0<n$, luego $B_r^{pq}=B_\infty(E_2^{pq})$ para $r<j_0-q+1$.
    Por otra parte, $\mathrm H^n(F_{\mathrm I}^p(K^{\bullet\bullet}))
    =\mathrm H^n(K^{\bullet\bullet})$ para $p+j_0<n-1$.
\end{enumerate}
\end{proof}

\begin{env}[11.3.4]
\label{0.11.3.4-fr}
Supongamos que el bicomplejo $K^{\bullet\bullet}=(K^{ij})$ sea tal que $K^{ij}=0$ para $i<0$ o $j<0$. Se sabe que se puede entonces definir para todo $p\in\mathbb Z$ un «homomorfismo de borde» canónico
\[
\label{0.11.3.4.1-fr}
  'E_2^{p0}(K^{\bullet\bullet})\longrightarrow
  \mathrm H^p(K^{\bullet\bullet})
  \tag{11.3.4.1}
\]
(M, XV, 6). Recordemos brevemente que esto se debe, por una parte, a que se tiene en la sucesión espectral $'E(K^{\bullet\bullet})$, $Z_r^{p0}=Z_{\mathrm I}^p(Z_{\mathrm{II}}^0(K^{\bullet\bullet}))$ para $2\leq r\leq+\infty$, y por otra al hecho de que $\mathrm H^p(F_{\mathrm I}^{p+1}(K^{\bullet\bullet}))=0$, de modo que el isomorfismo $\beta^{p0}:'E_\infty^{p0}\xrightarrow{\sim}
\mathrm H^p(F_{\mathrm I}^p)/\mathrm H^p(F_{\mathrm I}^{p+1})$ da un homomorfismo $'E_\infty^{p0}\to\mathrm H^p(F_{\mathrm I}^p(K^{\bullet\bullet}))
\to\mathrm H^p(K^{\bullet\bullet})$; la igualdad de todos los $Z_r^{p0}$ permite entonces definir homomorfismos canónicos $'E_r^{p0}\to'E_s^{p0}$ para $r\leq s$, y en particular un homomorfismo $'E_2^{p0}\to'E_\infty^{p0}$, de donde, por composición, el homomorfismo de borde $'E_2^{p0}\to\mathrm H^p(K^{\bullet\bullet})$; además, se verifica inmediatamente que, a la clase mód. $B_2^{p0}$ de un elemento $z\in Z_{\mathrm{II}}^0(K^{\bullet\bullet})\subset K^{p0}$ tal que $d'z=0$, el homomorfismo de borde así definido hace corresponder en $'E_\infty^{p0}$ la clase de $z$ mód. $B_\infty^{p0}$, y después a esta última la clase de cohomología de $z$ en $\mathrm H^p(K^{\bullet\bullet})$. Se ve, pues, finalmente que el \emph{homomorfismo de borde} (\hyperref[0.11.3.4.1-fr]{11.3.4.1}) \emph{proviene, por paso a la cohomología, de la inyección canónica} $Z_{\mathrm{II}}^0(K^{\bullet\bullet})\to K^{\bullet\bullet}$ (donde $K^{\bullet\bullet}$ se considera como complejo simple). Naturalmente, se interpreta del mismo modo el homomorfismo de borde
\[
\label{0.11.3.4.2-fr}
  ''E_2^{p0}(K^{\bullet\bullet})\longrightarrow
  \mathrm H^p(K^{\bullet\bullet})
  \tag{11.3.4.2}
\]
como proveniente de la inyección canónica $Z_{\mathrm I}^0(K^{\bullet\bullet})\to K^{\bullet\bullet}$.
\end{env}

\begin{env}[11.3.5]
\label{0.11.3.5-fr}
Sea ahora $K_{\bullet\bullet}=(K_{ij})$ un bicomplejo de
$\boldsymbol{C}$ cuyos dos operadores de derivación son de grado $-1$. Escribiremos
entonces $K_{i,\bullet}$ (resp. $K_{\bullet,j}$) para el complejo simple
$(K_{ij})_{j\in\mathbb Z}$ (resp. $(K_{ij})_{i\in\mathbb Z}$),
$\mathrm H_p^{\mathrm{II}}(K_{i,\bullet})$ (resp.
$\mathrm H_p^{\mathrm I}(K_{\bullet,j})$) para su objeto de homología de grado $p$,
$\mathrm H_p^{\mathrm{II}}(K_{\bullet\bullet})$ (resp.
$\mathrm H_p^{\mathrm I}(K_{\bullet\bullet})$) para el complejo
$(\mathrm H_p^{\mathrm{II}}(K_{i,\bullet}))_{i\in\mathbb Z}$ (resp.
$(\mathrm H_p^{\mathrm I}(K_{\bullet,j}))_{j\in\mathbb Z}$),
$\mathrm H_q^{\mathrm I}(\mathrm H_p^{\mathrm{II}}(K_{\bullet\bullet}))$
(resp.
$\mathrm H_q^{\mathrm{II}}(\mathrm H_p^{\mathrm I}(K_{\bullet\bullet}))$)
para su objeto de homología de grado $q$; notaciones análogas para los objetos de ciclos
y objetos de bordes; finalmente, $\mathrm H_n(K_{\bullet\bullet})$ designará
(cuando exista) el objeto de homología de grado $n$ del complejo simple (con
operador de derivación de grado $-1$) definido por $K_{\bullet\bullet}$.

Sea $K^{\prime\bullet\bullet}=(K^{\prime ij})$ con
$K^{\prime ij}=K_{-i,-j}$ el bicomplejo con operadores de derivación de grados
$+1$ asociado a $K_{\bullet\bullet}$; por definición, las
\emph{sucesiones espectrales} de $K_{\bullet\bullet}$ son las de
$K^{\prime\bullet\bullet}$, que se escriben $'E(K_{\bullet\bullet})$ y
$''E(K_{\bullet\bullet})$, donde sin embargo se cambian las notaciones, poniendo
\[
  'E_{pq}^r(K_{\bullet\bullet})
  ='E_r^{-p,-q}(K^{\prime\bullet\bullet}),\qquad
  ''E_{pq}^r(K_{\bullet\bullet})
  =''E_r^{-p,-q}(K^{\prime\bullet\bullet})
\]
para $2\leq r\leq\infty$. Con estas notaciones, se tiene
\[
  'E_{pq}^2(K_{\bullet\bullet})
  =\mathrm H_p^{\mathrm I}(\mathrm H_q^{\mathrm{II}}(K_{\bullet\bullet})),
  \qquad
  ''E_{pq}^2(K_{\bullet\bullet})
  =\mathrm H_p^{\mathrm{II}}(\mathrm H_p^{\mathrm I}(K_{\bullet\bullet})).
\]
Para evitar errores de signo, será en general preferible, para las
relaciones entre estas sucesiones espectrales y su término límite, volver al
complejo $K^{\prime\bullet\bullet}$. Señalemos no obstante los criterios
correspondientes a (\hyperref[0.11.3.3-fr]{11.3.3}):
\end{env}

\begin{env}[11.3.6]
\label{0.11.3.6}
\label{0.11.3.6-fr}
Las sucesiones espectrales $'E(K_{\bullet\bullet})$ y $''E(K_{\bullet\bullet})$ son \emph{birregulares} en los casos siguientes:
\begin{enumerate}
  \item[\textit{a)}] Existen $i_0$ y $j_0$ tales que $K_{ij}=0$ para $i>i_0$ o $j>j_0$ (resp. para $i<i_0$
    \oldpage[0\textsubscript{III}]{32}
    o $j<j_0$);
  \item[\textit{b)}] Existen $i_0$ e $i_1$ tales que $K_{ij}=0$ para $i<i_0$ y $i>i_1$;
  \item[\textit{c)}] Existen $j_0$ y $j_1$ tales que $K_{ij}=0$ para $j<j_0$ y $j>j_1$.
\end{enumerate}

La sucesión $'E(K_{\bullet\bullet})$ es \emph{regular} si existe $i_0$ tal que $K_{ij}=0$ para $i<i_0$, o si existe $j_0$ tal que $K_{ij}=0$ para $j>j_0$.

La sucesión $''E(K_{\bullet\bullet})$ es \emph{regular} si existe $i_0$ tal que $K_{ij}=0$ para $i>i_0$, o si existe $j_0$ tal que $K_{ij}=0$ para $j<j_0$.
\end{env}


\subsection{Hipercohomología de un funtor con respecto a un complejo $K^\bullet$}
\label{subsection:0.11.4}
\label{subsection:0.11.4-fr}

\begin{env}[11.4.1]
\label{0.11.4.1}
\label{0.11.4.1-fr}
Sea $\boldsymbol{C}$ una categoría abeliana.
Recordemos que se ha definido una \emph{resolución por la derecha} (o \emph{resolución cohomológica}) de un objeto $A$ de $\boldsymbol{C}$ como un complejo de objetos de $\boldsymbol{C}$, cuyo operador diferencial es de grado $+1$,
\[
  0\longrightarrow L^0\longrightarrow L^1\longrightarrow L^2
  \ldots
\]
provisto de un morfismo $\varepsilon:A\to L^0$, llamado la \emph{aumentación} de la resolución, que puede considerarse como un morfismo de complejos
\[
\xymatrix{
  0\ar[r] & A\ar[r]\ar[d]_{\varepsilon} & 0\ar[r]\ar[d]
    & 0\ar[r]\ar[d] & \ldots\\
  0\ar[r] & L^0\ar[r] & L^1\ar[r] & L^2\ar[r] & \ldots
})
\]
tal que la sucesión
\[
  0\longrightarrow A\xrightarrow{\varepsilon}L^0\longrightarrow L^1
  \longrightarrow\ldots
\]
sea \emph{exacta}.
Análogamente, una \emph{resolución por la izquierda} (o \emph{resolución homológica}) de $A$ es un complejo $0\leftarrow L_0\leftarrow L_1\leftarrow\ldots$ de objetos de $\boldsymbol{C}$, cuyo operador diferencial es de grado~$-1$, provisto de una aumentación $\varepsilon:L_0\to A$, de modo que la sucesión
\[
  0\longleftarrow A\xleftarrow{\varepsilon}L_0\longleftarrow L_1
  \longleftarrow\ldots
\]
sea \emph{exacta}.

Cuando la resolución por la derecha $(L_i)_{i\geq0}$ de un objeto $A$ es tal que $L_i=0$ para $i\geq n+1$, se dice que esta resolución es \emph{de longitud $\leq n$}.
Se define análogamente una resolución por la izquierda de longitud $\leq n$.
Una resolución que es de longitud $\leq n$ para algún entero $n$ se dice \emph{finita}.

Una resolución de $A$ se llama \emph{proyectiva} (resp. \emph{inyectiva}) si los objetos de $\boldsymbol{C}$ distintos de $A$ que la componen son \emph{proyectivos} (resp. \emph{inyectivos}).
Cuando $\boldsymbol{C}$ es la categoría de los módulos (a izquierda, por ejemplo) sobre un anillo, se dice que una resolución de $A$ es \emph{plana} (resp. \emph{libre}) cuando los módulos distintos de $A$ que la componen son \emph{planos} (resp. \emph{libres}).
\end{env}

\begin{env}[11.4.2]
\label{0.11.4.2}
\label{0.11.4.2-fr}
Sea $K^\bullet=(K^i)_{i\in\mathbb Z}$ un complejo de objetos de $\boldsymbol{C}$, cuyo operador diferencial es de grado~$+1$.
Se define una \emph{resolución de Cartan--Eilenberg por la derecha} de $K^\bullet$ como un par formado por un bicomplejo $L^{\bullet\bullet}=(L^{ij})$ con operadores diferenciales de grado~$+1$ y tal que $L^{ij}=0$ para $j<0$, y un morfismo de complejos simples $\varepsilon:K^\bullet\to L^{\bullet,0}$, de modo que se cumplan las condiciones siguientes:
\oldpage[0\textsubscript{III}]{33}
\begin{enumerate}
  \item[(i)] Para cada índice $i$, las sucesiones
    \[
    \begin{aligned}
      0\longrightarrow{}&K^i\xrightarrow{\varepsilon}L^{i0}
        \longrightarrow L^{i1}\longrightarrow\ldots\\
      0\longrightarrow{}&\mathrm B^i(K^\bullet)\xrightarrow{\varepsilon}
        \mathrm B_{\mathrm I}^i(L^{\bullet,0})
        \longrightarrow\mathrm B_{\mathrm I}^i(L^{\bullet,1})
        \longrightarrow\ldots\\
      0\longrightarrow{}&\mathrm Z^i(K^\bullet)\xrightarrow{\varepsilon}
        \mathrm Z_{\mathrm I}^i(L^{\bullet,0})
        \longrightarrow\mathrm Z_{\mathrm I}^i(L^{\bullet,1})
        \longrightarrow\ldots\\
      0\longrightarrow{}&\mathrm H^i(K^\bullet)\xrightarrow{\varepsilon}
        \mathrm H_{\mathrm I}^i(L^{\bullet,0})
        \longrightarrow\mathrm H_{\mathrm I}^i(L^{\bullet,1})
        \longrightarrow\ldots
    \end{aligned}
    \]
    son exactas, es decir, $(L^{i,\bullet})$, $(\mathrm B_{\mathrm I}^i(L^{\bullet\bullet}))$, $(\mathrm Z_{\mathrm I}^i(L^{\bullet\bullet}))$ y
$(\mathrm H_{\mathrm I}^i(L^{\bullet\bullet}))$ son, respectivamente, \emph{resoluciones} de $K^i$, $\mathrm B^i(K^\bullet)$, $\mathrm Z^i(K^\bullet)$ y $\mathrm H^i(K^\bullet)$.
  \item[(ii)] Para cada $j$, el complejo simple $L^{\bullet,j}$ es \emph{escindido} o, en otras palabras, las sucesiones exactas
    \[
      \label{0.11.4.2.1-fr}
      0\longrightarrow\mathrm B_{\mathrm I}^i(L^{\bullet,j})
      \longrightarrow\mathrm Z_{\mathrm I}^i(L^{\bullet,j})
      \longrightarrow\mathrm H_{\mathrm I}^i(L^{\bullet,j})
      \longrightarrow0
      \tag{11.4.2.1}
    \]
    \[
      \label{0.11.4.2.2-fr}
      0\longrightarrow\mathrm Z_{\mathrm I}^i(L^{\bullet,j})
      \longrightarrow L^{ij}
      \longrightarrow\mathrm B_{\mathrm I}^{i+1}(L^{\bullet,j})
      \longrightarrow0
      \tag{11.4.2.2}
    \]
    son \emph{escindidas}.
\end{enumerate}

Se demuestra (M,~XVII,~1.2) que si todo objeto de $\boldsymbol{C}$ es subobjeto de un objeto inyectivo, entonces todo complejo $K^\bullet$ de $\boldsymbol{C}$ admite una resolución de Cartan--Eilenberg \emph{inyectiva}, es decir, formada por objetos inyectivos $L^{ij}$ (la condición (ii) anterior implica entonces que los $\mathrm B_{\mathrm I}^i(L^{\bullet,j})$, $\mathrm Z_{\mathrm I}^i(L^{\bullet,j})$ y
$\mathrm H_{\mathrm I}^i(L^{\bullet,j})$ están también formados por objetos inyectivos).
Además, para todo morfismo $f:K^\bullet\to K'^\bullet$ de complejos de $\boldsymbol{C}$,
toda resolución de Cartan--Eilenberg $L^{\bullet\bullet}$ de $K^\bullet$ y toda resolución de Cartan--Eilenberg \emph{inyectiva} $L'^{\bullet\bullet}$ de $K'^\bullet$, existe un morfismo de bicomplejos $F:L^{\bullet\bullet}\to L'^{\bullet\bullet}$ compatible con $f$ y las aumentaciones, y si $f$ y $g$ son morfismos homótopos de $K^\bullet$ en $K'^\bullet$, entonces los morfismos correspondientes de $L^{\bullet\bullet}$ en $L'^{\bullet\bullet}$ son también homótopos \emph{(loc. cit.)}.

Cuando $K^\bullet$ está \emph{acotado inferiormente} (resp. \emph{superiormente}), puede tomarse $L^{\bullet\bullet}$ tal que $L^{ij}=0$ para $i<i_0$ (resp. $i>i_0$) si $K^i=0$ para $i<i_0$ (resp. $i>i_0$) (M,~XVII,~1.3).

Supongamos, por otra parte, que existe un entero $n$ tal que todo objeto de $\boldsymbol{C}$ admite una \emph{resolución inyectiva de longitud $\leq n$};
entonces puede suponerse que $L^{ij}=0$ para $j>n$ (M,~XVII,~1.4).
\end{env}


\begin{env}[11.4.3]
\label{0.11.4.3}
\label{0.11.4.3-fr}
Sea ahora $T$ un \emph{funtor covariante aditivo} de $\boldsymbol{C}$ en una categoría abeliana $\boldsymbol{C}'$.
Dado un complejo $K^\bullet$ de $\boldsymbol{C}$ y una resolución de Cartan--Eilenberg \emph{inyectiva} $L^{\bullet\bullet}$ de $K^\bullet$, supongamos que existe el complejo (simple) definido por el bicomplejo $T(L^{\bullet\bullet})$ existe
(cf. \hyperref[0.11.3.1-fr]{11.3.1}); entonces las dos sucesiones espectrales
$'E(T(L^{\bullet\bullet}))$ y $''E(T(L^{\bullet\bullet}))$ de este bicomplejo se llaman las \emph{sucesiones espectrales de hipercohomología} de $T$ con respecto a $K^\bullet$; en virtud de (\hyperref[0.11.4.2-fr]{11.4.2}) y (\hyperref[0.11.3.2-fr]{11.3.2}), efectivamente solo dependen de $K^\bullet$ y no de la elección de la resolución de Cartan--Eilenberg inyectiva $L^{\bullet\bullet}$;
además, dependen \emph{funtorialmente} de $K^\bullet$.
Tienen el mismo término límite $\mathrm H^\bullet(T(L^{\bullet\bullet}))$, llamado también la \emph{hipercohomología} de $T$ con respecto a $K^\bullet$, y denotado por $\mathbf R^\bullet T(K^\bullet)$.
Se demuestra que los términos $E_2$ de las dos sucesiones espectrales anteriores vienen dados por
\[
  \label{0.11.4.3.1-fr}
  'E_2^{pq}=\mathrm H^{\boldsymbol{p}}(\mathrm R^qT(K^\bullet))
  \tag{11.4.3.1}
\]
\[
  \label{0.11.4.3.2-fr}
  ''E_2^{pq}=\mathrm R^{\boldsymbol{p}}T(\mathrm H^q(K^\bullet))
  \tag{11.4.3.2}
\]
\oldpage[0\textsubscript{III}]{34}
donde $\mathrm R^pT$ denota, como de costumbre, el \emph{funtor derivado} de grado $p$ de $T$ para $p\in\mathbb Z$;
$\mathrm R^qT(K^\bullet)$ denota el complejo $(\mathrm R^qT(K^i))_{i\in\mathbb Z}$.
\phantomsection\label{0.11.4.4}
\label{0.11.4.4-fr}
Salvo mención explícita en contrario, supondremos en lo sucesivo que todo objeto de $\boldsymbol{C}$ es subobjeto de un objeto inyectivo de $\boldsymbol{C}$, de modo que existen resoluciones de Cartan--Eilenberg inyectivas para todo complejo de $\boldsymbol{C}$.
Como $L^{ij}=0$ para $j<0$, los criterios de \sref{0.11.3.3} muestran que las \emph{dos} sucesiones espectrales de hipercohomología de $T$ con respecto a $K^\bullet$ existen y
son \emph{birregulares} en cada uno de los dos casos siguientes:
\begin{enumerate}
  \item[$1^\circ$] $K^\bullet$ está acotado inferiormente;
  \item[$2^\circ$] Todo objeto de $\boldsymbol{C}$ admite una resolución inyectiva
de longitud a lo sumo igual a un entero $n$ independiente del objeto considerado.
\end{enumerate}
En efecto, puede suponerse, en el primer caso, por \sref{0.11.4.2}, que existe $i_0$ tal que $L^{ij}=0$ para $i<i_0$, y, en el segundo caso, que existe $j_1$ tal que $L^{ij}=0$ para $j>j_1$;
en cada uno de estos dos casos, está además claro que, para cada $n$ dado, solo existe un número finito de pares $(i,j)$ tales que $L^{ij}\neq0$ e $i+j=n$, lo que demuestra nuestras aserciones.

Si suponemos que los límites inductivos filtrantes existen en $\boldsymbol{C}'$ y son exactos (lo que implica, en particular, la existencia de sumas directas infinitas en $\boldsymbol{C}'$), entonces existe el complejo definido por el bicomplejo $T(L^{\bullet\bullet})$, y el criterio \sref{0.11.3.3} muestra que la sucesión $'E(T(L^{\bullet\bullet}))$ es siempre \emph{regular}.

Si $K^\bullet$ es un complejo tal que todos los $K^i$ son nulos salvo un único $K^{i_0}$, entonces
$\mathbf R^nT(K^\bullet)$ es isomorfo a $\mathrm R^{n-i_0}T(K^{i_0})$, como se deduce inmediatamente de las definiciones tomando una resolución de Cartan--Eilenberg $L^{\bullet\bullet}$ tal que $L^{ij}=0$ para $i\neq i_0$.

Si $K^\bullet$ y $K'^\bullet$ son dos complejos de $\boldsymbol{C}$, y $f$ y $g$ son morfismos homótopos de $K^\bullet$ en $K'^\bullet$, entonces los morfismos $\mathbf R^\bullet T(K^\bullet)\to
\mathbf R^\bullet T(K'^\bullet)$ inducidos por $f$ y $g$ son idénticos, y lo mismo ocurre con los morfismos de las sucesiones espectrales de hipercohomología.
\end{env}

\begin{proposition}[11.4.5]
\label{0.11.4.5-fr}
Supongamos que en $\boldsymbol{C}'$ los límites inductivos filtrantes existen y son exactos. Si $\mathrm R^nT(K^i)=0$ para todo $n>0$ y todo $i\in\mathbb Z$, se tienen isomorfismos funtoriales
\[
  \label{0.11.4.5.1-fr}
  \mathbf R^iT(K^\bullet)\xrightarrow{\sim}
  \mathrm H^i(T(K^\bullet))\qquad(i\in\mathbb Z).
  \tag{11.4.5.1}
\]
\end{proposition}

\begin{proof}
En efecto, los únicos términos $E_2$ no nulos de la primera sucesión espectral \sref{0.11.4.3.1} son entonces $'E_2^{p0}=\mathrm H^p(T(K^\bullet))$;
en otras palabras, esta sucesión es \emph{degenerada};
como es regular \sref{0.11.4.4}, la conclusión se deduce de \sref{0.11.1.6}.
\end{proof}

\begin{env}[11.4.6]
\label{0.11.4.6}
\label{0.11.4.6-fr}
Consideremos ahora, por ejemplo, un \emph{bifuntor} covariante $(M,N)\to T(M,N)$ de $\boldsymbol{C}\times\boldsymbol{C}'$ en $\boldsymbol{C}''$, donde $\boldsymbol{C}$, $\boldsymbol{C}'$ y $\boldsymbol{C}''$ son tres categorías abelianas;
se supone, por simplicidad, que $T$ es aditivo en cada uno de sus argumentos y, además, que todo objeto de $\boldsymbol{C}$ y todo objeto de $\boldsymbol{C}'$ son subobjetos de un objeto inyectivo,
y que los límites inductivos filtrantes existen en $\boldsymbol{C}''$ y son exactos.
Se define entonces la \emph{hipercohomología} de $T$ con respecto a los complejos $K^\bullet$ y $K'^\bullet$ de $\boldsymbol{C}$ y $\boldsymbol{C}'$, respectivamente, con operadores diferenciales de grado~$+1$, reemplazando $K^\bullet$ (resp. $K'^\bullet$) por una resolución de Cartan--Eilenberg inyectiva $L^{\bullet\bullet}$ (resp. $L'^{\bullet\bullet}$);
entonces $T(L^{\bullet\bullet},L'^{\bullet\bullet})$ es un cuadricomplejo de $\boldsymbol{C}''$, que se considera como un \emph{bicomplejo} de $\boldsymbol{C}''$, siendo los grados de $T(L^{ij},L'^{hk})$ los enteros $i+h$ y $j+k$.
La \emph{hipercohomología} de $T$ con respecto a $K^\bullet$ y $K'^\bullet$ es, por definición, la cohomología $\mathrm H^\bullet(T(L^{\bullet\bullet},L'^{\bullet\bullet}))$ de este bicomplejo (en otras palabras, la del complejo simple asociado) y se denota por $\mathbf R^\bullet T(K^\bullet,K'^\bullet)$;
\oldpage[0\textsubscript{III}]{35}
es el término límite de dos sucesiones espectrales cuyos términos $E_2$ vienen dados por
\[
  \label{0.11.4.6.1-fr}
  'E_2^{pq}=\mathrm H^p(\mathrm R^qT(K^\bullet,K'^\bullet))
  \tag{11.4.6.1}
\]
\[
  \label{0.11.4.6.2-fr}
  ''E_2^{pq}=\sum_{q'+q''=q}
  \mathrm R^pT(\mathrm H^{q'}(K^\bullet),\mathrm H^{q''}(K'^\bullet))
  \tag{11.4.6.2}
\]
(cf.~M,~XVII,~2).

Aquí $\mathrm R^qT(K^\bullet,K'^\bullet)$ es el bicomplejo $(\mathrm R^qT(K^i,K'^j))_{(i,j)\in\mathbb Z\times\mathbb Z}$, y el segundo miembro de \sref{0.11.4.6.1} es su cohomología cuando se lo considera como complejo simple.

Además, la primera sucesión espectral es siempre regular, y las dos sucesiones espectrales son birregulares cuando existe $n$ tal que todo objeto de $\boldsymbol{C}$ y todo objeto de $\boldsymbol{C}'$ admiten una resolución inyectiva de longitud $\leq n$, o cuando $K^\bullet$ \emph{y} $K'^\bullet$ están acotados inferiormente;
en este último caso, puede omitirse además la hipótesis de que existan los límites inductivos en $\boldsymbol{C}$ y $\boldsymbol{C}'$.

Si $K_1^\bullet$ y $K_1'^\bullet$ son otros dos complejos de $\boldsymbol{C}$ y $\boldsymbol{C}'$, respectivamente, $f$ y $g$ son morfismos \emph{homótopos} de $K^\bullet$ en $K_1^\bullet$,
y $f'$ y $g'$ son morfismos \emph{homótopos} de $K'^\bullet$ en $K_1'^\bullet$, entonces los morfismos $\mathbf R^\bullet T(K^\bullet,K'_\bullet)\to
\mathbf R^\bullet T(K_1^\bullet,K_1'^\bullet)$ inducidos por $f$ y $f'$, por una parte, y por $g$ y $g'$, por otra, son idénticos, y lo mismo ocurre con los morfismos de las sucesiones espectrales de hipercohomología.

Se generaliza fácilmente a todo multifuntor covariante aditivo.
\end{env}

\begin{proposition}[11.4.7]
\label{0.11.4.7}
\label{0.11.4.7-fr}
Supongamos que, para todo objeto inyectivo $I$ de $\boldsymbol{C}$ \textup{(resp.} $I'$ de $\boldsymbol{C}'$), $A'\leadsto T(I,A')$ \textup{(resp.} $A\leadsto T(A,I')$) es un funtor exacto.
Entonces, con la notación de \sref{0.11.4.6}, se tiene isomorfismos canónicos
\[
  \label{0.11.4.7.1-fr}
  \mathbf R^\bullet T(K^\bullet,K'^\bullet)
  \xrightarrow{\sim}\mathrm H^\bullet(T(L^{\bullet\bullet},K'^\bullet))
  \xrightarrow{\sim}\mathrm H^\bullet(T(K^\bullet,L'^{\bullet\bullet}))
  \tag{11.4.7.1}
\]
donde los dos últimos términos son la cohomología de los complejos simples definidos por los tricomplejos $T(L^{\bullet\bullet},K'^\bullet)$ y $T(K^\bullet,L'^{\bullet\bullet})$, respectivamente.
\end{proposition}

\begin{proof}
Definamos, por ejemplo, el primero de estos isomorfismos.
El cuadricomplejo $T(L^{\bullet\bullet},L'^{\bullet\bullet}))$ puede considerarse como un bicomplejo, donde los grados de $T(L^{ij},L'^{hk})$ son los enteros $i+j+h$ y $k$.
Como, para cada $h$, $L'^{h,\bullet}$ es una \emph{resolución} de $K'^h$, se tiene, para este bicomplejo, en virtud de las hipótesis sobre $T$, $\mathrm H_{\mathrm{II}}^q(T(L^{\bullet\bullet},L'^{\bullet\bullet}))=0$ para $q\neq0$, y $\mathrm H_{\mathrm{II}}^0(T(L^{\bullet\bullet},L'^{\bullet\bullet}))
=T(L^{\bullet\bullet},K'^\bullet)$;
la primera sucesión espectral de este bicomplejo es, pues, \emph{degenerada};
como $L'^{hk}=0$ para $k<0$, esta sucesión es también regular \sref{0.11.3.3}, y la conclusión se deduce entonces de \sref{0.11.1.6}.
\end{proof}

Se tienen resultados análogos para un multifuntor covariante de un número cualquiera $n$ de argumentos: en el cálculo de la hipercohomología no es necesario reemplazar \emph{todos} los complejos
por una resolución de Cartan--Eilenberg, sino únicamente $n-1$ de ellos, siempre que, al fijar $n-1$ argumentos arbitrarios tomándolos como objetos \emph{inyectivos}, el funtor covariante en el argumento restante sea \emph{exacto}.


\subsection{Paso al límite inductivo en la hipercohomología}
\label{subsection:0.11.5}
\label{subsection:0.11.5-fr}

\begin{lemma}[11.5.1]
\label{0.11.5.1-fr}
Sea $K^\bullet=(K^i)_{i\in\mathbb Z}$ un complejo de $\boldsymbol{C}$, y para todo entero $r\in\mathbb Z$, sea $K_{(r)}^\bullet$ el complejo tal que $K_{(r)}^i=0$ para $i<r$, $K_{(r)}^i=K^i$ para $i\geq r$. Sea $T$ un funtor covariante aditivo
\oldpage[0\textsubscript{III}]{36}
de $\boldsymbol{C}$ en $\boldsymbol{C}'$, que conmuta con los límites inductivos (se supone que los límites inductivos filtrantes existen y son exactos en $\boldsymbol{C}$ y $\boldsymbol{C}'$). Entonces $\mathbf R^\bullet T(K^\bullet)$ es isomorfo al límite inductivo $\varinjlim\mathbf R^\bullet T(K_{(r)}^\bullet)$ cuando $r$ tiende a $-\infty$.
\end{lemma}

La construcción de una resolución inyectiva de Cartan-Eilenberg de $K^\bullet$ se hace eligiendo \emph{arbitrariamente}, para cada $i$, una resolución inyectiva $(X_{\mathrm B}^{ij})_{j\geq0}$ de $\mathrm B^i(K^\bullet)$ y una resolución inyectiva $(X_{\mathrm H}^{ij})_{j\geq0}$ de $\mathrm H^i(K^\bullet)$; hecho esto, el método de construcción muestra que la resolución inyectiva $(L^{ij})_{j\geq0}$ de $K^i$ y los operadores de derivación $L^{i,j}\to L^{i+1,j}$ \emph{solo dependen de las resoluciones} $(X_{\mathrm B}^{i,\bullet})$, $(X_{\mathrm H}^{i,\bullet})$ \emph{y} $(X_{\mathrm B}^{i+1,\bullet})$ (M, XVII, 1.2). Ahora bien, está claro que se tiene $\mathrm B^i(K_{(r)}^\bullet)=\mathrm B^i(K^\bullet)$ y $\mathrm H^i(K_{(r)}^\bullet)=\mathrm H^i(K^\bullet)$ para $i\geq r+1$. Por otra parte, para cada $r$, se tiene una inyección canónica $\varphi_{r-1,r}^\bullet:K_{(r)}^\bullet\to K_{(r-1)}^\bullet$, siendo $\varphi_{r-1,r}^i$ la identidad para $i\neq r-1$. La observación precedente muestra que, si $L^{\bullet\bullet}=(L^{ij})$ es una resolución inyectiva de Cartan-Eilenberg de $K^\bullet$, se puede para cada $r$ definir una resolución inyectiva de Cartan-Eilenberg $L_{(r)}^{\bullet\bullet}=(L_{(r)}^{ij})$ de $K_{(r)}^\bullet$ tal que $L_{(r)}^{ij}=0$ para $i<r$ y $L_{(r)}^{ij}=L^{ij}$ para $i\geq r+1$. Por otra parte, se puede definir un morfismo de bicomplejos $\Phi_{r-1,r}^{\bullet\bullet}:L_{(r)}^{\bullet\bullet}\to
L_{(r-1)}^{\bullet\bullet}$ correspondiente a $\varphi_{r-1,r}^\bullet$, y el método de definición de este morfismo (\emph{loc. cit.}) muestra además que puede construirse de modo que $\Phi_{r-1,r}^{ij}$ sea la identidad para $i\neq r$ y $i\neq r-1$. Se ha definido así un sistema inductivo $(L_{(r)}^{\bullet\bullet})$ de bicomplejos de $\boldsymbol{C}$, del que $L^{\bullet\bullet}$ es evidentemente el límite inductivo cuando $r$ tiende a $-\infty$; por la permutabilidad de las sumas directas y de los límites inductivos, el complejo simple asociado a $L^{\bullet\bullet}$ es también límite inductivo del complejo simple asociado a $L_{(r)}^{\bullet\bullet}$. Como $T$ conmuta por hipótesis con los límites inductivos y lo mismo ocurre con la cohomología (por la exactitud del funtor $\varinjlim$), se tiene en efecto $\mathrm H^\bullet(T(L^{\bullet\bullet}))=\varinjlim\mathrm H^\bullet(T(L_{(r)}^{\bullet\bullet}))$ salvo isomorfismo.

El lema (\hyperref[0.11.5.1-fr]{11.5.1}) permite extender a complejos $K^\bullet$ cualesquiera, por paso al límite inductivo, propiedades válidas para los complejos \emph{acotados inferiormente}. Como primer ejemplo, demostraremos:

\begin{proposition}[11.5.2]
\label{0.11.5.2}
\label{0.11.5.2-fr}
Bajo las hipótesis de \sref{0.11.5.1} relativas a $\boldsymbol{C}$, $\boldsymbol{C}'$ y $T$, el funtor $\mathbf R^\bullet T(K^\bullet)$ es un funtor cohomológico sobre la categoría abeliana de los complejos de $\boldsymbol{C}$.
\end{proposition}

\begin{proof}
Mostremos primero que basta tratar los complejos \emph{acotados inferiormente}.
Dada una sucesión exacta
$0\longrightarrow K'^\bullet\longrightarrow K^\bullet\longrightarrow K''^\bullet\longrightarrow0$
de complejos, para todo $r$ induce una sucesión exacta
$0\longrightarrow K_{(r)}'^\bullet\longrightarrow K_{(r)}^\bullet\longrightarrow K_{(r)}''^\bullet\longrightarrow0$,
Por hipótesis, esta última da una sucesión exacta
\[
  \ldots\longrightarrow\mathbf R^nT(K_{(r)}'^\bullet)
  \longrightarrow\mathbf R^nT(K_{(r)}^\bullet)
  \longrightarrow\mathbf R^nT(K_{(r)}''^\bullet)
  \xrightarrow{\partial}\mathbf R^{n+1}T(K_{(r)}'^\bullet)
  \longrightarrow\ldots
\]
Estas sucesiones exactas forman un sistema inductivo.
El Lema~\sref{0.11.5.1} y la exactitud de $\varinjlim$ dan, por tanto, una sucesión exacta
\[
  \ldots\longrightarrow\mathbf R^nT(K'^\bullet)
  \longrightarrow\mathbf R^nT(K^\bullet)
  \longrightarrow\mathbf R^nT(K''^\bullet)
  \xrightarrow{\partial}\mathbf R^{n+1}T(K'^\bullet)
  \longrightarrow\ldots
\]
\end{proof}

Para tratar los complejos acotados inferiormente, basta considerar aquellos para los cuales $K^i=0$ cuando $i<0$; forman evidentemente una \emph{categoría abeliana} $\boldsymbol{K}$.

\begin{lemma}[11.5.2.1]
\label{0.11.5.2.1}
\label{0.11.5.2.1-fr}
En $\boldsymbol{K}$, sea $\mathfrak J$ la clase de los complejos $Q^\bullet=(Q^i)_{i\geq0}$ que tienen las
\oldpage[0\textsubscript{III}]{37}
propiedades siguientes: $1^\circ$ Todo $Q^i$ es un objeto inyectivo de $\boldsymbol{C}$; $2^\circ$ Para todo $i\geq0$, $\mathrm Z^i(Q^\bullet)=\mathrm B^i(Q^\bullet)$, y $\mathrm Z^i(Q^\bullet)$ es un sumando directo de $Q^i$.
Entonces:
\begin{enumerate}
  \item[{\rm(i)}] todo $Q^\bullet\in\mathfrak J$ es un objeto inyectivo de $\boldsymbol{K}$;
  \item[{\rm(ii)}] todo objeto de $\boldsymbol{K}$ es isomorfo a un subcomplejo de un complejo perteneciente a $\mathfrak J$.
\end{enumerate}
\end{lemma}

\begin{proof}
{\rm(i)}
Sea $A^\bullet=(A^i)$ un objeto de $\boldsymbol{K}$, sea $A'^\bullet=(A'^i)$ un subobjeto de $A^\bullet$, sea $Q^\bullet=(Q^i)$ perteneciente a $\mathfrak J$, y supongamos dado un morfismo
$f=(f^i):A'^\bullet\to Q^\bullet$, que se trata de prolongar a un morfismo $g=(g^i):A^\bullet\to Q^\bullet$.
Por simplicidad, utilizamos el lenguaje de la categoría de módulos (cf.~[27]).

Identifiquemos $Q$ con $\mathrm B^i(Q^\bullet)\oplus\mathrm B^{i+1}(Q^\bullet)$.
Procedemos por inducción sobre $i$.
Supongamos que $g^j$ ha sido definido para $j<i$, compatiblemente con los diferenciales $d^j:A^j\to A^{j+1}$ y $d'^j:Q^j\to Q^{j+1}$ para $j<i-1$ y tales además que: $1^\circ$ $g^{i-1}(\mathrm Z^{i-1}(A^\bullet))\subset
\mathrm Z^{i-1}(Q^\bullet)$; $2^\circ$ Si se pone $C^j=(d^j)^{-1}(A'^{j+1})$ para todo $j$, entonces $d'^{i-1}\circ g^{i-1}$ coincide con $f^i\circ d^{i-1}$ sobre $C^{i-1}$.
Componiendo $f^i:A'^i\to Q^i$ con las proyecciones, se obtienen dos morfismos $f'^i:A'^i\to\mathrm B^i(Q^\bullet)$ y $f''^i:A'^i\to\mathrm B^{i+1}(Q^\bullet)$. Como $d'^{i-1}\circ g^{i-1}$ envía $A^{i-1}$ en $\mathrm B^i(Q^\bullet)$ y se anula sobre $\mathrm Z^{i-1}(A^\bullet)$, define un morfismo
$h^i:\mathrm B^i(A^\bullet)\to\mathrm B^i(Q^\bullet)$, y puesto que $d'^{i-1}\circ g^{i-1}$ coincide con $f^i\circ d^{i-1}$ sobre $C^{i-1}$, $h^i$ coincide con $f'_i$ sobre $\mathrm B^i(A^\bullet)\cap A'^i$.
Como $\mathrm B^i(Q^\bullet)$ es un sumando directo de $Q^i$, es inyectivo; existe, por tanto, un morfismo
$g'^i:A^i\to\mathrm B^i(Q^\bullet)$
que coincide con $h^i$ sobre $\mathrm B^i(A^\bullet)$ y con $f'^i$ sobre $A'^i$.

Por otra parte, consideremos el morfismo
$f''^{i+1}\circ d^i:C^i\to\mathrm B^{i+1}(Q^\bullet)$,
que se anula sobre $\mathrm Z^i(A^\bullet)$.
Como $\mathrm B^{i+1}(Q^\bullet)$ es inyectivo, existe un morfismo
$g''^i:A^i\to\mathrm B^{i+1}(Q^\bullet)$
que coincide con $f''^{i+1}\circ d^i$ sobre $C^i$ y con $0$ en $\mathrm Z^i(A^\bullet)$.
Tomando $g^i=g'^i+g''^i$ puede continuarse la inducción.

\smallskip
\noindent{\rm(ii)}
Para sumergir $A^\bullet=(A^i)$ en un complejo perteneciente a $\mathfrak J$, se toma para cada $i\geq1$ un objeto inyectivo $Q'^i$ de $\boldsymbol{C}$ y una inyección $f'^i:A^i\to Q'^i$. Pongamos entonces $Q^i=0$ para $i<0$, $Q^0=Q'^1$ y $Q^i=Q'^i\oplus Q'^{i+1}$ para $i\geq1$, con el operador de derivación evidente. Pongamos
$f^i=f'^i+(f'^{i+1}\circ d^i)$ para todo $i$ (con $f'^i=0$ para $i\leq0$);
Es inmediato que $f^i$ es inyectivo para todo $i$, y que los $f^i$ conmutan con los diferenciales.
\end{proof}

\begin{corollary}[11.5.2.2]
\label{0.11.5.2.2}
\label{0.11.5.2.2-fr}
Todo objeto $K^\bullet$ de $\boldsymbol{K}$ admite una resolución por la derecha formada por objetos de $\mathfrak J$.
Si $L^{\bullet\bullet}$ es una tal resolución, entonces para toda resolución $L'^{\bullet\bullet}$ de $K^\bullet$ formada por objetos de $\boldsymbol{K}$ existe un morfismo de bicomplejos
$F:L'^{\bullet\bullet}\to L^{\bullet\bullet}$ compatible con las aumentaciones, y dos tales morfismos $F$, $F'$ son homótopos.
\end{corollary}

\begin{proof}
Esto no es sino (M, V, 1.1 \emph{a)}) aplicado a la categoría abeliana
$\boldsymbol{K}$.
\end{proof}

\begin{env}[11.5.2.3]
\label{0.11.5.2.3-fr}
Establecidos estos preliminares, consideremos una resolución de Cartan-Eilenberg inyectiva $L'^{\bullet\bullet}$ una resolución de Cartan--Eilenberg inyectiva de $K^\bullet$ y sea $L^{\bullet\bullet}$ una resolución de $K^\bullet$ formada por objetos de $\mathfrak J$.
Mostremos que existe un isomorfismo
$\mathrm H^\bullet(T(L'^{\bullet\bullet}))\xrightarrow{\sim}\mathrm H^\bullet(T(L^{\bullet\bullet}))$.
El Corolario~\sref{0.11.5.2.2} da un morfismo de bicomplejos
$T(L'^{\bullet\bullet})\to T(L^{\bullet\bullet})$,
y, por tanto, un morfismo
$'E(T(L'^{\bullet\bullet}))\to{}'E(T(L^{\bullet\bullet}))$
entre sus primeras sucesiones espectrales.
Por \sref{0.11.3.3}, estas sucesiones espectrales son \emph{regulares}, de modo que por \sref{0.11.1.5} basta comprobar que este morfismo es un isomorfismo sobre los términos $E_2$, o también que
$\mathrm H_{\mathrm{II}}^q(T(L'^{i,\bullet}))$ es igual a
$\mathrm R^qT(K^i)$; como $L'^{i,\bullet}$ es una resolución por la derecha de $K^i$, esto reduce la cuestión al lema siguiente.
\end{env}

\oldpage[0\textsubscript{III}]{38}
\begin{lemma}[11.5.2.4]
\label{0.11.5.2.4}
\label{0.11.5.2.4-fr}
Sea $L^\bullet=(L^i)_{i\in\mathbb Z}$ una resolución por la derecha de un objeto $A$ de $\boldsymbol{C}$ tal que $\mathrm R^nT(L^i)=0$ para todo $i\in\mathbb Z$ y todo $n>0$.
Entonces
$\mathrm H^\bullet T(L^\bullet)=\mathrm R^\bullet T(A)$.
\end{lemma}

\begin{proof}
Este es un caso particular de (T, 2.5.2).
\end{proof}

\begin{env}[11.5.2.5]
\label{0.11.5.2.5}
\label{0.11.5.2.5-fr}
La demostración de \sref{0.11.5.2} se deduce ahora inmediatamente, puesto que $L'^{\bullet\bullet}$ es una resolución inyectiva de $K^\bullet$ \emph{en la categoría abeliana $\boldsymbol{K}$}.
En otras palabras, el funtor
$K^\bullet\leadsto\mathrm H^\bullet(T(L'^{\bullet\bullet}))$
es precisamente el funtor cohomológico \emph{derivado por la derecha} de $T$ en la categoría $\boldsymbol{K}$ (T, 2.3).
\end{env}

\begin{proposition}[11.5.3]
\label{0.11.5.3}
\label{0.11.5.3-fr}
Bajo las hipótesis de \sref{0.11.5.1} relativas a $\boldsymbol{C}$, $\boldsymbol{C}'$ y $T$, sea $L^{\bullet\bullet}=(L^{ij})$ un bicomplejo tal que $L^{ij}=0$ para $j<0$ y tal que, para todo $i$, $L^{i,\bullet}$ es una resolución de $K^i$.
Supongamos además que $\mathrm R^nT(L^{ij})=0$ para todo par $(i,j)$ y todo $n>0$.
Entonces existe un isomorfismo funtorial
\[
  \label{0.11.5.3.1-fr}
  \mathbf R^\bullet T(K^\bullet)\xrightarrow{\sim}
  \mathrm H^\bullet(T(L^{\bullet\bullet}))).
  \tag{11.5.3.1}
\]
\end{proposition}

\begin{proof}
Sea $L_{(r)}^{\bullet\bullet}=(L_{(r)}^{ij})$ el bicomplejo tal que $L_{(r)}^{ij}=0$ para $i<r$, $L_{(r)}^{ij}=L^{ij}$ para $i\geq r$; es inmediato que $L^{\bullet\bullet}$ es el límite inductivo de los $L_{(r)}^{\bullet\bullet}$ cuando $r$ tiende a $-\infty$.
Por la hipótesis sobre $T$ y por \sref{0.11.5.1}, basta, por tanto, demostrar la proposición cuando $K^\bullet$ está acotado inferiormente, por ejemplo cuando $K^i=0$ para $i<0$, con $L^{ij}=0$ para $i<0$.
Sea $L'^{\bullet\bullet}=(L'^{ij})$ una resolución por la derecha de $K^\bullet$ formada por objetos de $\mathfrak J$, como en \sref{0.11.5.2.2}.
Existe un morfismo de bicomplejos
$L^{\bullet\bullet}\to L'^{\bullet\bullet}$
compatible con las aumentaciones y, por tanto, un morfismo
$'E(T(L^{\bullet\bullet}))\to{}'E(T(L'^{\bullet\bullet}))$
entre las primeras sucesiones espectrales.
El Lema~\sref{0.11.5.2.4} muestra, exactamente como en \sref{0.11.5.2.3}, que este morfismo es un isomorfismo, lo que demuestra la aserción.
\end{proof}

\begin{remark}[11.5.4]
\label{0.11.5.4}
\label{0.11.5.4-fr}
Los razonamientos anteriores muestran que las conclusiones de \sref{0.11.5.2} y \sref{0.11.5.3} se cumplen en la categoría de los complejos $K^\bullet$ \emph{acotados inferiormente} sin suponer que $T$ conmuta con los límites inductivos filtrantes.
Además, si se considera únicamente la categoría $\boldsymbol{K}$ de los complejos $K^\bullet$ tales que $K^i=0$ para $i<0$, la caracterización de $\mathbf R^\bullet T(K^\bullet)$ como el sistema de funtores derivados por la derecha de $T$ en $\boldsymbol{K}$ muestra que este funtor cohomológico es \emph{universal} (T, 2.3).

Otro caso en el que \sref{0.11.5.2} se cumple sin hipótesis adicional alguna sobre $T$ es aquel en que existe un entero $m>0$ tal que todo objeto de $\boldsymbol{C}$ admita una resolución inyectiva \emph{de longitud} $\leq m$. En efecto, en la demostración de (\hyperref[0.11.5.1-fr]{11.5.1}), todas las resoluciones inyectivas de objetos de $\boldsymbol{C}$ que intervienen pueden tomarse de longitud $\leq m$, de donde resulta inmediatamente que los términos de grado total $n$ en $T(L_{(r)}^{\bullet\bullet})$ son \emph{iguales a los de} $T(L^{\bullet\bullet})$ y en número \emph{finito}, en cuanto $r$ es suficientemente grande, lo que
implica que para todo $n$,
$\mathrm H^n(T(L^{\bullet\bullet}))=
\mathrm H^n(T(L_{(r)}^{\bullet\bullet}))$
para $r$ suficientemente grande.
Con la notación de \sref{0.11.5.2}, se tiene análogamente
$\mathbf R^nT(K_{(r)}^\bullet)=\mathbf R^nT(K^\bullet)$
para $r$ suficientemente grande (dependiendo de $n$), y análogamente para $K'^\bullet$ y $K''^\bullet$, lo que demuestra la conclusión.
Del mismo modo, \sref{0.11.5.3} se cumple sin condición adicional sobre $T$ cuando $\boldsymbol{C}$ satisface la hipótesis anterior y se toman las resoluciones $L^{i,\bullet}$ son de longitud $\leq m$.
\end{remark}

\begin{env}[11.5.5]
\label{0.11.5.5-fr}
El resultado de (\hyperref[0.11.5.2-fr]{11.5.2}) se generaliza a los multifuntores covariantes. Consideremos por ejemplo la situación de (\hyperref[0.11.4.6-fr]{11.4.6}), donde se supone que en $\boldsymbol{C}$, $\boldsymbol{C}'$ y $\boldsymbol{C}''$
\oldpage[0\textsubscript{III}]{39}
los límites inductivos filtrantes existen y son exactos, y que $T$ conmuta con los límites inductivos. Entonces $\mathbf R^\bullet T(K^\bullet,K'^\bullet)$ es un bifuntor \emph{cohomológico} en cada uno de los complejos $K^\bullet$, $K'^\bullet$; para verlo, se reduce, como en (\hyperref[0.11.5.2-fr]{11.5.2}), al caso en que $K^\bullet$ y $K'^\bullet$ están acotados inferiormente; tomando entonces para $K^\bullet$ y $K'^\bullet$ resoluciones inyectivas del tipo descrito en (\hyperref[0.11.5.2.2-fr]{11.5.2.2}), se reduce a la propiedad general demostrada en (M, V, 4.1).
\end{env}

\begin{env}[11.5.6]
\label{0.11.5.6-fr}
Del mismo modo, los resultados de (\hyperref[0.11.4.7-fr]{11.4.7}) y (\hyperref[0.11.5.3-fr]{11.5.3}) se generalizan como sigue. Supongamos (bajo las hipótesis de (\hyperref[0.11.5.5-fr]{11.5.5})) que se tienen dos bicomplejos $L^{\bullet\bullet}=(L^{ij})$, $L'^{\bullet\bullet}=(L'^{ij})$ tales que $L^{ij}=0$ y $L'^{ij}=0$ para $j<0$, tales que, para todo $i$, $L^{i,\bullet}$ es una resolución de $K^i$ y $L'^{i,\bullet}$ es una resolución de $K'^i$, y tales que
$\mathrm R^nT(L^{ij},L'^{hk})=0$
para $n>0$ y todo sistema de índices $(i,j,h,k)$.
Entonces existe un isomorfismo, funtorial en $K^\bullet$ y $K'^\bullet$,
\[
  \label{0.11.5.6.1-fr}
  \mathbf R^\bullet T(K^\bullet,K'^\bullet)\xrightarrow{\sim}
  \mathrm H^\bullet(T(L^{\bullet\bullet},L'^{\bullet\bullet}))
  \tag{11.5.6.1}
\]
Esto se demuestra como en \sref{0.11.5.3}, reduciéndose al caso en que $K^\bullet$ y $K'^\bullet$ están acotados inferiormente.

Supongamos además que, para todo par $(i,j)$ y todo par $(h,k)$, los funtores $A\leadsto T(A,L'^{hk})$ y $A'\leadsto T(L^{ij},A')$ sean exactos en $\boldsymbol{C}$ y $\boldsymbol{C}'$, respectivamente.
Entonces existen también isomorfismos funtoriales
\[
  \label{0.11.5.6.2-fr}
  \mathbf R^\bullet T(K^\bullet,K'^\bullet)
  \sim\mathrm H^\bullet(T(L^{\bullet\bullet},K'^\bullet))
  \xrightarrow{\sim}\mathrm H^\bullet(T(K^\bullet,L'^{\bullet\bullet})).
  \tag{11.5.6.2}
\]
La demostración es análoga a la de \sref{0.11.4.7}.
\end{env}

\begin{env}[11.5.7]
\label{0.11.5.7}
\label{0.11.5.7-fr}
Los resultados de \sref{0.11.5.5} y \sref{0.11.5.6} se cumplen también sin suponer que $T$ conmuta con los límites inductivos, con tal de restringirse a los complejos $K^\bullet$, $K'^\bullet$ acotados inferiormente, o de suponer que todo objeto de $\boldsymbol{C}$ (resp. de $\boldsymbol{C}'$) admite una resolución inyectiva de longitud acotada y, en \sref{0.11.5.6}, que los segundos grados de $L^{\bullet\bullet}$ y $L'^{\bullet\bullet}$ están acotados superiormente.
\end{env}

\subsection{Hiperhomología de un funtor con respecto a un complejo $K_\bullet$}
\label{subsection:0.11.6}
\label{subsection:0.11.6-fr}

\begin{env}[11.6.1]
\label{0.11.6.1}
\label{0.11.6.1-fr}
Sea $\boldsymbol{C}$ una categoría abeliana y sea $K_\bullet=(K_i)_{i\in\mathbb Z}$ un complejo de objetos de $\boldsymbol{C}$ cuyo diferencial es de grado $-1$.
Una \emph{resolución de Cartan--Eilenberg por la izquierda} de $K_\bullet$ consiste en un bicomplejo $L_{\bullet\bullet}=(L_{ij})$, cuyos diferenciales son de grado $-1$ y para el cual $L_{ij}=0$ cuando $j<0$, junto con un morfismo de complejos simples
$\varepsilon:L_{\bullet,0}\to K_\bullet$,
tal que se satisfacen las condiciones obtenidas de las de \sref{0.11.4.2} ``invirtiendo las flechas''.
Si todo objeto de $\boldsymbol{C}$ es cociente de un objeto proyectivo, entonces todo complejo $K_\bullet$ de $\boldsymbol{C}$ admite una \emph{resolución de Cartan--Eilenberg proyectiva}, es decir, formada por objetos proyectivos $L_{ij}$; posee propiedades funtoriales análogas a las de \sref{0.11.4.2}.
Además, si $K_\bullet$ está acotado inferiormente (resp. superiormente), puede exigirse $L_{ij}=0$ para $i<i_0$ (resp. $i>i_0$) siempre que $K_i=0$ para $i<i_0$ (resp. $i>i_0$).
Si todo objeto de $\boldsymbol{C}$ admite una resolución proyectiva de longitud $\leq n$, se puede suponer que $L_{ij}=0$ para $j>n$.
\end{env}

\begin{env}[11.6.2]
\label{0.11.6.2}
\label{0.11.6.2-fr}
\oldpage[0\textsubscript{III}]{40}
Supongamos que $T$ es un \emph{funtor covariante aditivo} de $\boldsymbol{C}$ en una categoría abeliana $\boldsymbol{C}'$.
La \emph{hiperhomología} $\mathbf{L}_\bullet T(K_\bullet)$ y las \emph{sucesiones espectrales de hiperhomología} de $T$ con respecto a un complejo $K_\bullet$ de $\boldsymbol{C}$ (cuando existen) se definen a partir de \sref{0.11.4.3}, de nuevo ``invirtiendo las flechas''.
Los términos $E^2$ de las dos sucesiones espectrales así obtenidas son
\[
\label{0.11.6.2.1-fr}
\tag{11.6.2.1}
{}'E^2_{pq}=H_p(L_qT(K_\bullet))
\]
y
\[
\label{0.11.6.2.2-fr}
\tag{11.6.2.2}
{}''E^2_{pq}=L_pT(H_q(K_\bullet))
\]
donde $L_pT$ denota el funtor derivado de grado $p$ de $T$ para $p\geq 0$, y $0$ para $p<0$, mientras que $L_qT(K_\bullet)$ denota el complejo $(L_qT(K_i))_{i\in\mathbb Z}$.

Las propiedades de la hiperhomología no se deducen todas por simple «inversión de las flechas» de las de la hipercohomología (a menos que se hagan hipótesis suplementarias del tipo AB~5*) de T, 1.5 sobre la categoría $\boldsymbol{C}'$) a causa de las condiciones de \emph{regularidad} de las dos sucesiones espectrales precedentes, a las que esta vez hay que aplicar los criterios de (\hyperref[0.11.3.4-fr]{11.3.4}). Estos últimos muestran que, cuando se supone que en $\boldsymbol{C}'$ los límites inductivos filtrantes existen y son exactos, entonces el complejo definido por el bicomplejo $T(L_{\bullet\bullet})$ existe, y la \emph{segunda} sucesión espectral $\,{}''E(T(L_{\bullet\bullet}))$ es esta vez \emph{regular}. Si se supone, o bien que $K_\bullet$ está acotado inferiormente, o bien que existe un entero $n$ tal que todo objeto de $\boldsymbol{C}$ admita una resolución proyectiva de longitud $\leq n$, entonces las \emph{dos} sucesiones espectrales de hiperhomología existen (sin hipótesis sobre $\boldsymbol{C}'$) y son \emph{birregulares}.
\end{env}

\begin{proposition}[11.6.3]
\label{0.11.6.3}
\label{0.11.6.3-fr}
Sean $\boldsymbol{C}$ y $\boldsymbol{C}'$ dos categorías abelianas, y sea $T$ un funtor covariante aditivo de $\boldsymbol{C}$ en $\boldsymbol{C}'$.
Entonces:
\begin{enumerate}
  \item[{\rm(i)}] {La hiperhomología $\mathbf{L}_\bullet T(K_\bullet)$ es un funtor homológico sobre la categoría abeliana de los complejos de $\boldsymbol{C}$ acotados inferiormente.}

  \item[{\rm(ii)}] {Sea $K_\bullet$ un complejo de $\boldsymbol{C}$ acotado inferiormente.
  Si $L_nT(K_i)=0$ para todo $n>0$ y todo $i\in\mathbb Z$, existen isomorfismos funtoriales}
  \[
\label{0.11.6.3.1-fr}
\tag{11.6.3.1}
\mathbf{L}_iT(K_\bullet)\xrightarrow{\sim} H_i(T(K_\bullet))\qquad(i\in\mathbb Z).
\]

  \item[{\rm(iii)}] {Sea $K_\bullet$ un complejo de $\boldsymbol{C}$ acotado inferiormente.
  Sea $L_{\bullet\bullet}=(L_{ij})$ un bicomplejo tal que $L_{ij}=0$ para $j<0$ y, para todo $i$, $L_{i,\bullet}$ es una resolución de $K_i$.
  Supongamos además que $L_nT(L_{ij})=0$ para todo par $(i,j)$ y todo $n>0$.
  Entonces existe un isomorfismo funtorial}
  \[
\label{0.11.6.3.2-fr}
\tag{11.6.3.2}
\mathbf{L}_\bullet T(K_\bullet)\xrightarrow{\sim} H_\bullet(T(L_{\bullet\bullet}))
\]
\end{enumerate}

\upshape Las demostraciones proceden como las de \sref{0.11.5.2}, \sref{0.11.4.5} y \sref{0.11.5.3} en el caso de complejos acotados inferiormente.
Dejamos al lector los detalles de estos razonamientos.
\end{proposition}

\begin{env}[11.6.4]
\label{0.11.6.4}
\label{0.11.6.4-fr}
Se tienen resultados enteramente análogos para los \emph{multifuntores} covariantes aditivos en cada argumento.
Por ejemplo, para un bifuntor $T$, las dos sucesiones espectrales de hiperhomología tienen términos $E^2$
\oldpage[0\textsubscript{III}]{41}
\[
\label{0.11.6.4.1-fr}
\tag{11.6.4.1}
{}'E^2_{pq}=H_p(L_qT(K_\bullet,K'_\bullet))
\]
y
\[
\label{0.11.6.4.2-fr}
\tag{11.6.4.2}
{}''E^2_{pq}=
\sum_{q'+q''=q}L_p(H_{q'}(K_\bullet),H_{q''}(K'_\bullet)).
\]
También aquí es la \emph{segunda} sucesión espectral la que es regular.
Las dos sucesiones son birregulares cuando $K_\bullet$ y $K'_\bullet$ están acotados inferiormente, o cuando los objetos de las categorías abelianas consideradas admiten resoluciones proyectivas de una longitud fija.

Además:
\end{env}

\begin{proposition}[11.6.5]
\label{0.11.6.5}
\label{0.11.6.5-fr}
Sean $\boldsymbol{C}$, $\boldsymbol{C}'$ y $\boldsymbol{C}''$ tres categorías abelianas, y sea $T$ un bifuntor covariante biaditivo de $\boldsymbol{C}\times\boldsymbol{C}'$ en $\boldsymbol{C}''$.
\begin{enumerate}
  \item[{\rm(i)}] {$\mathbf{L}_\bullet T(K_\bullet,K'_\bullet)$ es un bifuntor homológico en cada uno de los complejos acotados inferiormente $K_\bullet$, $K'_\bullet$, formados respectivamente por objetos de $\boldsymbol{C}$ y de $\boldsymbol{C}'$.}

  \item[{\rm(ii)}] {Supongamos que $K_\bullet$ y $K'_\bullet$ están acotados inferiormente.
  Sean $L_{\bullet\bullet}=(L_{ij})$ y $L'_{\bullet\bullet}=(L'_{ij})$ dos bicomplejos tales que $L_{ij}=0$ y $L'_{ij}=0$ para $j<0$ y tales que, para todo $i$, $L_{i,\bullet}$ es una resolución de $K_i$ y $L'_{i,\bullet}$ es una resolución de $K'_i$, y finalmente que $L_nT(L_{ij},L'_{hk})=0$ para todo $n>0$ y todo sistema de índices $(i,j,h,k)$.
  Entonces existe un isomorfismo funtorial}
  \[
\label{0.11.6.5.1-fr}
\tag{11.6.5.1}
\mathbf{L}_\bullet T(K_\bullet,K'_\bullet)
\xrightarrow{\sim} H_\bullet(T(L_{\bullet\bullet},L'_{\bullet\bullet})).
\]

  \item[{\rm(iii)}] {Supongamos además que, para todo par $(i,j)$ y todo par $(h,k)$, los funtores $A\leadsto T(A,L'_{hk})$ y $A'\leadsto T(L_{ij},A')$ sean exactos en $\boldsymbol{C}$ y $\boldsymbol{C}'$, respectivamente.
  Entonces existen isomorfismos funtoriales}
  \[
\label{0.11.6.5.2-fr}
\tag{11.6.5.2}
\mathbf{L}_\bullet T(K_\bullet,K'_\bullet)
\xrightarrow{\sim} H_\bullet(T(L_{\bullet\bullet},K'_\bullet))
\xrightarrow{\sim} H_\bullet(T(K_\bullet,L'_{\bullet\bullet})).
\]
\end{enumerate}

\upshape Las demostraciones son análogas a las de \sref{0.11.5.5} y \sref{0.11.5.6}.
\end{proposition}

\subsection{Hiperhomología de un funtor con respecto a un bicomplejo $K_{\bullet\bullet}$}
\label{subsection:0.11.7}
\label{subsection:0.11.7-fr}

\begin{env}[11.7.1]
\label{0.11.7.1}
\label{0.11.7.1-fr}
Sea $\boldsymbol{C}$ una categoría abeliana en la que todo objeto es cociente de un
objeto proyectivo.
Consideremos un bicomplejo $K_{\bullet\bullet}=(K_{ij})$ de objetos de $\boldsymbol{C}$ cuyos
dos grados están \emph{acotados inferiormente}.
Se puede limitar siempre al caso en que $K_{ij}=0$ para $i<0$ o
$j<0$, y es lo que haremos en adelante. Se puede considerar $K_{\bullet\bullet}$ como un complejo ordinario formado por los
objetos
$K_{i,\bullet}=(K_{ij})_{j\geq 0}$
de la categoría abeliana $\boldsymbol{K}$ de los complejos no negativos de objetos de
$\boldsymbol{C}$.
Del Lema~\sref{0.11.5.2.1} (o, más precisamente, del lema
``dual'' obtenido ``invirtiendo las flechas'') y de (M,~V,~2.2) se deduce
que $K_{\bullet\bullet}$ admite una \emph{resolución de Cartan--Eilenberg
proyectiva} en la categoría $\boldsymbol{K}$.
Una tal resolución es un \emph{tricomplejo}
$M_{\bullet\bullet\bullet}=(M_{ijk})$ en $\boldsymbol{C}$, con grados todos $\geq 0$, formado por objetos proyectivos, tal que para todo $i$, $M_{i,\bullet\bullet}$, $\mathrm B_i^{\mathrm I}(M_{\bullet\bullet\bullet})$, $\mathrm Z_i^{\mathrm I}(M_{\bullet\bullet\bullet})$, $H_i^{\mathrm I}(M_{\bullet\bullet\bullet})$ constituyen \emph{resoluciones proyectivas} de $K_{i,\bullet}$, $\mathrm B_i^{\mathrm I}(K_{\bullet\bullet})$, $\mathrm Z_i^{\mathrm I}(K_{\bullet\bullet})$, $H_i^{\mathrm I}(K_{\bullet\bullet})$ respectivamente en la categoría $\boldsymbol{K}$; en particular, para todo par $(i,j)$, $M_{ij,\bullet}$ es una \emph{resolución
proyectiva} de $K_{ij}$ en $\boldsymbol{C}$.
\end{env}

\begin{proposition}[11.7.2]
\label{0.11.7.2}
\label{0.11.7.2-fr}
Sea $T$ un funtor covariante aditivo de $\boldsymbol{C}$ en una categoría
abeliana $\boldsymbol{C}'$.
Con la notación de \sref{0.11.7.1}, la homología
$H_\bullet(T(M_{\bullet\bullet\bullet}))$
del complejo simple asociado al tricomplejo
$T(M_{\bullet\bullet\bullet})$ es canónicamente isomorfa a la
hiperhomología $\mathbf{L}_\bullet T(K_{\bullet\bullet})$ del
\oldpage[0\textsubscript{III}]{42}
complejo simple asociado a $K_{\bullet\bullet}$
\sref{0.11.6.2}.
Es también el término límite común de seis sucesiones espectrales birregulares denotadas
$\mathop{}^{(t)}E$, donde $t=a,b,a',b',c$ o $d$, cuyos
términos $E^2$ están dados por las fórmulas
\[
\begin{aligned}
{}^{(a)}E^2_{pq}&=L_pT(H_q(K_{\bullet\bullet}))&
{}^{(b)}E^2_{pq}&=H_p(\mathbf{L}_q^{\mathrm{II}}T(K_{\bullet\bullet}))\\
{}^{(a')}E^2_{pq}&=\mathbf{L}_pT(H_q^{\mathrm I}(K_{\bullet\bullet}))&
{}^{(b')}E^2_{pq}&=H_p(L_qT(K_{\bullet\bullet}))\\
{}^{(c)}E^2_{pq}&=\mathbf{L}_pT(H_q^{\mathrm{II}}(K_{\bullet\bullet}))&
{}^{(d)}E^2_{pq}&=H_p(\mathbf{L}_q^{\mathrm I}T(K_{\bullet\bullet}))
\end{aligned}
\]
\end{proposition}

(Se recuerda que utilizamos la notación $F(A_\bullet)$ para el complejo cuyos objetos son los $F(A_i)$ para todo complejo $A_\bullet=(A_i)$; por ejemplo $\mathbf{L}_q^{\mathrm{II}}T(K_{\bullet\bullet})$ denota el complejo
$(\mathbf{L}_q^{\mathrm{II}}T(K_{i,\bullet}))_{i\geq0}$,
donde $\mathbf{L}_q^{\mathrm{II}}T(K_{i,\bullet})$ es la hiperhomología de índice
$q$ del funtor $T$ relativa al complejo simple $K_{i,\bullet}$.)

\begin{proof}
Designemos por $L_\bullet$ el complejo simple asociado a
$K_{\bullet\bullet}$, de modo que
$L_i=\bigoplus_{r+s=i}K_{rs}$,
y pongamos
$N_{ij}=\bigoplus_{r+s=i}M_{rsj}$;
Está claro que, para todo $i$, $N_{i,\bullet}$ es una \emph{resolución proyectiva}
de $L_i$ en $\boldsymbol{C}$.
De \sref{0.11.6.3} y \sref{0.11.6.4} se deduce que existe un
isomorfismo funtorial
$
\mathbf{L}_\bullet T(L_\bullet)\xrightarrow{\sim} H_\bullet(T(N_{\bullet\bullet}))$;
Como el complejo simple asociado al bicomplejo
$T(N_{\bullet\bullet})$ es también el complejo simple asociado al
tricomplejo $T(M_{\bullet\bullet\bullet})$, esto demuestra la primera aserción.

Además, $\mathbf{L}_\bullet T(L_\bullet)$ es el término límite de las dos
sucesiones espectrales de hiperhomología \sref{0.11.6.2} de $T$ relativas al
complejo simple $L_\bullet$.
Estas son precisamente $\mathop{}^{(b')}E$ y $\mathop{}^{(a)}E$,
respectivamente.

Consideremos ahora $M_{\bullet\bullet\bullet}$ como un bicomplejo
$U_{\bullet\bullet}$, donde
$U_{ij}=\bigoplus_{r+s=j}M_{irs}$;
La homología $H_\bullet(T(M_{\bullet\bullet\bullet}))$ es también la
término límite de las dos sucesiones espectrales del bicomplejo
$T(U_{\bullet\bullet})$.
Para todo $i$, el bicomplejo $M_{i,\bullet\bullet}$ satisface las
hipótesis de \sref{0.11.6.3} relativas al complejo simple
$K_{i,\bullet}$.
Por consiguiente,
$H_q^{\mathrm{II}}(T(U_{i,\bullet}))=
\mathbf{L}_q^{\mathrm{II}}T(K_{i,\bullet}))$,
y la primera sucesión espectral de $T(U_{\bullet\bullet})$ es precisamente
$\mathop{}^{(b)}E$.

Por otra parte, para todo $r$, $M_{\bullet,r,\bullet}$ es una
resolución de Cartan--Eilenberg del complejo simple $K_{\bullet,r}$.
El cálculo efectuado en (M,~XV,~2) muestra que
$H_q^{\mathrm I}(T(M_{\bullet,rs}))=T(H_q^{\mathrm I}(M_{\bullet,rs}))$,
y, por tanto,
$H_q^{\mathrm I}(T(U_{\bullet,j}))=\bigoplus_{r+s=j}T(H_q^{\mathrm I}(M_{\bullet,rs}))$;
En otras palabras, el complejo simple
$H_q^{\mathrm I}(T(U_{\bullet\bullet}))$ es el complejo simple
asociado al bicomplejo
$T(H_q^{\mathrm I}(M_{\bullet\bullet\bullet})$.
Para todo $q$, $H_q^{\mathrm I}(M_{\bullet\bullet\bullet})$ es una resolución \emph{proyectiva} del complejo simple $H_q^{\mathrm I}(K_{\bullet\bullet})$ en la categoría $\boldsymbol{K}$; aplicando (\hyperref[0.11.6.3-fr]{11.6.3}), se ve que se tiene
\[
H_p^{\mathrm{II}}
\bigl(T(H_q^{\mathrm I}(M_{\bullet\bullet\bullet}))\bigr)
=\mathbf{L}_pT(H_q^{\mathrm I}(K_{\bullet\bullet})),
\]
y, por tanto, la sucesión espectral $\mathop{}^{(a')}E$.
Finalmente, las sucesiones espectrales $\mathop{}^{(c)}E$ y
$\mathop{}^{(d)}E$ se obtienen intercambiando los papeles de los
índices $i$ y $j$ en la definición del tricomplejo
$M_{\bullet\bullet\bullet}$ y aplicando los razonamientos anteriores al
tricomplejo resultante.
\end{proof}

Se llama a $\mathbf{L}_\bullet T(K_{\bullet\bullet})$ la \emph{hiperhomología} de
$T$ relativa al \emph{bicomplejo} $K_{\bullet\bullet}$.

\begin{remarks}[11.7.3]
\label{0.11.7.3}
\label{0.11.7.3-fr}
\emph{Observaciones}.
\begin{enumerate}
  \item[{\rm(i)}]
  De \sref{0.11.6.3} se deduce que
  $\mathbf{L}_\bullet T(K_{\bullet\bullet})$ es un \emph{funtor homológico} sobre la
  categoría de los bicomplejos $K_{\bullet\bullet}$ en $\boldsymbol{C}$ que están
  acotados inferiormente en cada grado.

\oldpage[0\textsubscript{III}]{43}
  \item[{\rm(ii)}]
  Sea $M_{\bullet\bullet\bullet}$ un tricomplejo en $\boldsymbol{C}$ tal que,
  para todo par $(r,s)$, $M_{rs,\bullet}$ sea una \emph{resolución} de $K_{rs}$,
  y tal que
  $L_nT(M_{ijk})=0$
  para toda terna $(i,j,k)$ y todo $n>0$.
  Entonces existe un isomorfismo
  $
\mathbf{L}_\bullet T(K_{\bullet\bullet})
\xrightarrow{\sim} H_\bullet(T(M_{\bullet\bullet\bullet}))$;
  En efecto, con la notación utilizada en la demostración de \sref{0.11.7.2},
  $N_{i,\bullet}$ es una resolución de $L_i$ tal que
  $L_nT(N_{ij})=0$ para todo $n>0$ y todo par $(i,j)$, por lo que
  basta aplicar \sref{0.11.6.3},~(iii).

  \item[{\rm(iii)}]
  Los resultados de \sref{0.11.7.2} se generalizan inmediatamente a los
  multifuntores covariantes.
  Por ejemplo, sea $\boldsymbol{C}'$ una segunda categoría abeliana en la que todo
  objeto sea cociente de un objeto proyectivo, sea
  $K'_{\bullet\bullet}$ un bicomplejo en $\boldsymbol{C}'$ cuyos dos grados estén
  acotados inferiormente, y sea $T$ un bifuntor covariante aditivo de
  $\boldsymbol{C}\times\boldsymbol{C}'$ en una categoría abeliana $\boldsymbol{C}''$.
  Si $L_\bullet$ y $L'_\bullet$ son los complejos simples asociados a
  $K_{\bullet\bullet}$ y $K'_{\bullet\bullet}$, respectivamente, se define
  la hiperhomología de $T$ relativa a los dos bicomplejos
  $K_{\bullet\bullet}$ y $K'_{\bullet\bullet}$ como
  $\mathbf{L}_\bullet T(L_\bullet,L'_\bullet)$;
  Aplicando \sref{0.11.6.4} y \sref{0.11.6.5}, se obtienen, como en
  \sref{0.11.7.2}, seis sucesiones espectrales que convergen a esta hiperhomología;
  dejamos al lector su formulación explícita.
\end{enumerate}
\end{remarks}

\subsection{Suplemento sobre la cohomología de los complejos simpliciales}
\label{subsection:0.11.8}
\label{subsection:0.11.8-fr}

\begin{env}[11.8.1]
\label{0.11.8.1-fr}
Sean $A$ un conjunto finito, $\Sigma(A)$ el conjunto de las sucesiones finitas $\sigma=(\alpha_0,\ldots,\alpha_h)$ de elementos de $A$ («símplices» de $A$); se pone $|\sigma|=\{\alpha_0,\ldots,\alpha_h\}$; se recuerda que el complejo de cadenas $\mathrm C_\bullet(A)$ es el grupo abeliano libre graduado engendrado por los elementos de $\Sigma(A)$, siendo $(\alpha_0,\ldots,\alpha_h)$ de \emph{grado} $h$, con una diferencial definida por $d(\alpha_0,\ldots,\alpha_h)=\sum_{j=0}^{h}(-1)^j(\alpha_0,\ldots,\widehat{\alpha}_j,\ldots,\alpha_h)$. El subgrupo $\mathrm D_\bullet(A)$ de $\mathrm C_\bullet(A)$ engendrado por las cadenas $\sigma=(\alpha_0,\ldots,\alpha_h)$ para las que dos de los $\alpha_i$ son iguales, y por las cadenas $\pi(\sigma)-\varepsilon_\pi.\sigma$, donde para toda permutación $\pi$, $\pi(\sigma)=(\alpha_{\pi(0)},\ldots,\alpha_{\pi(h)})$ y $\varepsilon_\pi$ es el signo de $\pi$, es un subcomplejo de $\mathrm C_\bullet(A)$ cuyos elementos se llaman cadenas \emph{degeneradas}; se pone $\mathrm L_\bullet(A)=\mathrm C_\bullet(A)/\mathrm D_\bullet(A)$, y se tiene un homomorfismo natural de complejos $p:\mathrm C_\bullet(A)\longrightarrow\mathrm L_\bullet(A)$. Por otra parte, se define un homomorfismo de complejos $j:\mathrm L_\bullet(A)\longrightarrow\mathrm C_\bullet(A)$ del modo siguiente: se ordena totalmente $A$; a la clase mód. $\mathrm D_\bullet(A)$ de un símplex $\sigma=(\alpha_0,\ldots,\alpha_h)$, se hace corresponder $0$ si dos de los $\alpha_i$ son iguales, y la sucesión $(\beta_0,\ldots,\beta_h)$ de los $\alpha_i$ ordenados \emph{en orden creciente} en caso contrario. Está claro que $p\circ
j$ es la \emph{identidad} de $\mathrm L_\bullet(A)$.
\end{env}

\begin{env}[11.8.2]
\label{0.11.8.2}
\label{0.11.8.2-fr}
Sea $B$ un segundo conjunto finito.
Si $d'$ y $d''$ son los diferenciales de
$\mathrm C_\bullet(A)$ y $\mathrm C_\bullet(B)$, el complejo producto
tensorial
$\mathrm C_\bullet(A)\otimes\mathrm C_\bullet(B)$
puede considerarse como el grupo abeliano libre engendrado por los elementos de
$\Sigma(A)\times\Sigma(B)$, con la diferencial
$d(\sigma,\tau)=(d'\sigma,\tau)+(-1)^{h+1}(\sigma,d''\tau)$ si $\operatorname{Card}|\sigma|=h+1$.

Los homomorfismos naturales
$\mathrm C_\bullet(A)\to\mathrm L_\bullet(A)$ y
$\mathrm C_\bullet(B)\to\mathrm L_\bullet(B)$ definen un homomorfismo
$p:\mathrm C_\bullet(A)\otimes\mathrm C_\bullet(B)\longrightarrow\mathrm L_\bullet(A)\otimes\mathrm L_\bullet(B)$,
siendo este último producto tensorial isomorfo a
$(\mathrm C_\bullet(A)\otimes\mathrm C_\bullet(B))/(\mathrm C_\bullet(A)\otimes\mathrm D_\bullet(B)+\mathrm D_\bullet(A)\otimes\mathrm C_\bullet(B))$.
Análogamente, utilizando los homomorfismos
$\mathrm L_\bullet(A)\to\mathrm C_\bullet(A)$ y
$\mathrm L_\bullet(B)\to\mathrm C_\bullet(B)$ definidos en
\sref{0.11.8.1}, con órdenes totales sobre $A$ y $B$, se obtiene un
homomorfismo
$j:\mathrm L_\bullet(A)\otimes\mathrm L_\bullet(B)\longrightarrow\mathrm C_\bullet(A)\otimes\mathrm C_\bullet(B)$
tal que $p\circ j$ es la identidad.
\end{env}

\oldpage[0\textsubscript{III}]{44}
Con esta notación:
\begin{proposition}[11.8.3]
\label{0.11.8.3}
\label{0.11.8.3-fr}
Existe una homotopía
$h:\mathrm C_\bullet(A)\otimes\mathrm C_\bullet(B)\longrightarrow\mathrm C_\bullet(A)\otimes\mathrm C_\bullet(B)$
tal que $h(\sigma,\tau)$ es una combinación lineal de pares de símplices
$(\sigma_i,\tau_i)$ que satisfacen
$|\sigma_i|\subset|\sigma|$ y $|\tau_i|\subset|\tau|$, y tal que,
si $f=j\circ p$, entonces
\[
\label{0.11.8.3.1-fr}
\tag{11.8.3.1}
f-1=h\circ d+d\circ h.
\]
\end{proposition}

\begin{proof}
Basta definir $h$ sobre cada par $(\sigma,\tau)$ de símplices, por
inducción sobre la suma de los grados de $\sigma$ y $\tau$, pues se puede tomar $h=0$ cuando esta suma es $0$. Sea
$\omega=f(\sigma,\tau)-(\sigma,\tau)-h(d(\sigma,\tau))$;
Por la hipótesis de inducción y la definición de $d$, se tiene
$\omega\in\mathrm C_\bullet(|\sigma|)\otimes\mathrm C_\bullet(|\tau|)$.
Además,
\[
d\omega=f(d(\sigma,\tau))-d(\sigma,\tau)-d(h(d(\sigma,\tau)))
=h(d(d(\sigma,\tau)))=0
\]
en virtud de (\hyperref[0.11.8.3.1-fr]{11.8.3.1}) y de la hipótesis de inducción. Ahora bien, se tiene $H_q(\mathrm C_\bullet(A))=0$ para $q>0$ (G, I, 3.7.4), por tanto también $H_q(\mathrm C_\bullet(A)\otimes\mathrm C_\bullet(B))=0$ para $q>0$, en virtud de la fórmula de Künneth (G, I, 5.5.2). Aplicando esta observación al reemplazar $A$ reemplazado por $|\sigma|$ y $B$ por
$|\tau|$, se ve que existe un elemento $\omega'$ de $\mathrm C_\bullet(|\sigma|)\otimes\mathrm C_\bullet(|\tau|)$ tal que $\omega=d\omega'$.
Tomar $h(\sigma,\tau)=\omega'$ verifica \eqref{0.11.8.3.1} para el par
$(\sigma,\tau)$ y permite continuar la inducción.
\end{proof}

\begin{env}[11.8.4]
\label{0.11.8.4}
\label{0.11.8.4-fr}
Con la notación de \sref{0.11.8.2}, escribamos
$(\sigma,\tau)\leq(\sigma',\tau')$
si $|\sigma|\subset|\sigma'|$ y $|\tau|\subset|\tau'|$.
Un \emph{sistema de coeficientes} $\mathscr S$ sobre
$\Sigma(A)\times\Sigma(B)$ consiste en una familia
$(\Gamma_{\sigma,\tau})$ de grupos abelianos, donde
$\Gamma_{\sigma,\tau}$ depende únicamente de los conjuntos $|\sigma|$ y
$|\tau|$, y una familia de homomorfismos $\Gamma_{\sigma,\tau}\longrightarrow\Gamma_{\sigma',\tau'}$, para $(\sigma,\tau)\leq(\sigma',\tau')$, que forman un \emph{sistema inductivo} para esta relación de preorden. Se define entonces un \emph{complejo de cocadenas} $\mathrm C^\bullet(A,B;\mathscr S)$ como el conjunto de las familias $\lambda=(\lambda(\sigma,\tau))$, donde $(\sigma,\tau)$ recorre $\Sigma(A)\times\Sigma(B)$, con $\lambda(\sigma,\tau)\in\Gamma_{\sigma,\tau}$ para todo par $(\sigma,\tau)$. La diferencial está dada del modo siguiente: si
$d(\sigma,\tau)=\sum_i\pm(\sigma_i,\tau_i)$,
entonces $|\sigma_i|\subset|\sigma|$ y $|\tau_i|\subset|\tau|$ para todo
$i$, y se pone
\[
d\lambda(\sigma,\tau)=\sum_i\pm\lambda_i(\sigma_i,\tau_i),
\]
donde $\lambda_i(\sigma_i,\tau_i)$ denota la imagen canónica de
$\lambda(\sigma_i,\tau_i)$ en $\Gamma_{\sigma,\tau}$. Diremos que una cocadena $\lambda\in\mathrm C^\bullet(A,B;\mathscr S)$ se llama
\emph{bialternada} si $\lambda(\sigma,\tau)=0$ siempre que $\sigma$
o $\tau$ tiene dos términos iguales, y si se tiene $\lambda(\pi(\sigma),\tau)=\varepsilon_\pi\lambda(\sigma,\tau)$ y $\lambda(\sigma,\pi'(\tau))=\varepsilon_{\pi'}\lambda(\sigma,\tau)$ para permutaciones cualesquiera $\pi$ y $\pi'$ de los índices.
Estas cocadenas engendran manifiestamente un \emph{subcomplejo}
$\mathrm L^\bullet(A,B;\mathscr S)$ de
$\mathrm C^\bullet(A,B;\mathscr S)$.
\end{env}

\begin{proposition}[11.8.5]
\label{0.11.8.5-fr}
{La inyección canónica
$\mathrm L^\bullet(A,B;\mathscr S)\longrightarrow\mathrm C^\bullet(A,B;\mathscr S)$
define un isomorfismo para la cohomología de estos dos complejos.}
\end{proposition}

\begin{proof}
Observemos que si $p$ y $j$ tienen el sentido definido en (\hyperref[0.11.8.2-fr]{11.8.2}), las aplicaciones ${}^tp:\lambda\to\lambda\circ p$ y ${}^tj:\lambda\to\lambda\circ j$ están definidas en $\mathrm L^\bullet(A,B;\mathscr S)$ y
$\mathrm C^\bullet(A,B;\mathscr S)$, respectivamente; la primera es precisamente
la inyección canónica.
Como ${}^tj\circ{}^tp$ es la identidad, solo queda demostrar
que ${}^tp\circ{}^tj$ es homótopa a la identidad.
Por \sref{0.11.8.3}, la aplicación
${}^th:\lambda\to\lambda\circ h$
está \emph{definida} sobre $\mathrm C^\bullet(A,B;\mathscr S)$, de modo que al trasponer
\eqref{0.11.8.3.1} se obtiene el resultado deseado.
\end{proof}

\begin{env}[11.8.6]
\label{0.11.8.6}
\label{0.11.8.6-fr}
\oldpage[0\textsubscript{III}]{45}
La Proposición~\sref{0.11.8.5} reduce el cálculo de la cohomología de
$\mathrm L^\bullet(A,B;\mathscr S)$ al de
$\mathrm C^\bullet(A,B;\mathscr S)$. Recordemos por otra parte que esta última es, en virtud del teorema de Eilenberg-Zilber (G, I, 3.10.2), canónicamente isomorfa a la cohomología del complejo de cocadenas definido como sigue: se forma el complejo de cadenas $\mathrm P_\bullet(A,B)$ de combinaciones lineales de pares
$(\sigma,\tau)\in\Sigma(A)\times\Sigma(B)$ para los cuales $\sigma$ y
$\tau$ tengan \emph{el mismo grado}; la diferencial de este complejo
viene dada por
$d:(\sigma,\tau)\leadsto\sum_{j,k}(-1)^{j+k}(\sigma_j,\tau_k)$ si
$d\sigma=\sum_j(-1)^j\sigma_j$ y $d\tau=\sum_k(-1)^k\tau_k$;
se tienen entonces dos homomorfismos canónicos de complejos
\[
f:\mathrm P_\bullet(A,B)\longrightarrow
\mathrm C_\bullet(A)\otimes\mathrm C_\bullet(B),\qquad
g:\mathrm C_\bullet(A)\otimes\mathrm C_\bullet(B)
\longrightarrow\mathrm P_\bullet(A,B),
\] y se demuestra (\emph{loc. cit.}) que hay \emph{homotopías} $h$ y $h'$ tales que
\[
f\circ g-1=d\circ h+h\circ d\qquad\text{y}\qquad
g\circ f-1=d\circ h'+h'\circ d.
\]
Además, se tiene $f(\sigma,\tau)\in\mathrm C_\bullet(|\sigma|)
\otimes\mathrm C_\bullet(|\tau|)$ y $g(\sigma,\tau)\in\mathrm P_\bullet(|\sigma|,|\tau|)$ y las homotopías $h$, $h'$ \emph{pueden tomarse de modo que} $h(\sigma,\tau)\in\mathrm C_\bullet(|\sigma|)
\otimes\mathrm C_\bullet(|\tau|)$ \emph{y} $h'(\sigma,\tau)\in\mathrm P_\bullet(|\sigma|,|\tau|)$. Este punto proviene de que la definición de $h(\sigma,\tau)$ y $h'(\sigma,\tau)$ pueden definirse por
\emph{inducción} sobre la suma de los grados de $\sigma$ y $\tau$, y de que los
$H^q(\mathrm C_\bullet(|\sigma|)\otimes\mathrm C_\bullet(|\tau|))$ y
$H^q(\mathrm P_\bullet(|\sigma|,|\tau|))$
son \emph{nulos} para $q>0$ (\emph{loc. cit.}); se razona entonces como en (\hyperref[0.11.8.3-fr]{11.8.3}) y resulta la conclusión. Se define entonces $\mathrm P^\bullet(A,B;\mathscr S)$ como el conjunto de las familias $\lambda=(\lambda(\sigma,\tau))$ donde $(\sigma,\tau)$ recorre los pares cuyos términos tienen el mismo grado, con $\lambda(\sigma,\tau)\in\Gamma_{\sigma,\tau}$, y como se tiene $d\sigma=\sum_j(-1)^j\sigma_j\in\mathrm C_\bullet(|\sigma|)$ y $d\tau=\sum_k(-1)^k\tau_k\in\mathrm C_\bullet(|\tau|)$,
\[
d\lambda(\sigma,\tau)=
\sum(-1)^{j+k}\lambda(\sigma_j,\tau_k)
\]
está definida y da el diferencial de
$\mathrm P^\bullet(A,B;\mathscr S)$. Dicho esto, las aplicaciones ${}^tf:\lambda\to\lambda\circ f$, ${}^tg:\lambda\to\lambda\circ g$, ${}^th:\lambda\to\lambda\circ h$ y ${}^th':\lambda\to\lambda\circ h'$ están todas \emph{definidas} en virtud de las observaciones precedentes;
${}^tf\circ{}^tg$ y ${}^tg\circ{}^tf$ son homótopas a la
identidad, dando el isomorfismo requerido entre la cohomología de
$\mathrm C^\bullet(A,B;\mathscr S)$ y la de
$\mathrm P^\bullet(A,B;\mathscr S)$.
\end{env}

\begin{env}[11.8.7]
\label{0.11.8.7}
\label{0.11.8.7-fr}
\emph{Observación}.
El razonamiento de \sref{0.11.8.3}, aplicado a
$\mathrm C_\bullet(A)$ y $\mathrm L_\bullet(A)$, muestra que si
$j$ y $p$ se definen como en \sref{0.11.8.1}, entonces
$f=j\circ p$ satisface de nuevo \eqref{0.11.8.3.1}, con
$|h(\sigma)|\subset|\sigma|$.
Como en \sref{0.11.8.5}, se obtiene un isomorfismo de la cohomología de
$\mathrm L^\bullet(A;\mathscr S)$ en la de
$\mathrm C^\bullet(A;\mathscr S)$, definiéndose los dos complejos de manera
evidente.
Este es el resultado cuya demostración se esboza en (G, I, 3.8.1).
\end{env}

\begin{env}[11.8.8]
\label{0.11.8.8}
\label{0.11.8.8-fr}
Volvamos ahora a la notación y las hipótesis de \sref{0.11.8.2}, y sea
$\mathscr S^\bullet=(\mathscr S^k)$
un \emph{complejo} \emph{de sistemas de coeficientes} sobre
$\Sigma(A)\times\Sigma(B)$.
Para cada $(\sigma,\tau)$, los grupos $\Gamma^k_{\sigma,\tau}$
forman, por tanto, un complejo de grupos abelianos $(k\in\mathbb Z)$, y se tienen los diagramas conmutativos
\[
\xymatrix{
\Gamma^k_{\sigma,\tau}\ar[r]\ar[d] &
\Gamma^{k+1}_{\sigma,\tau}\ar[d]\\
\Gamma^k_{\sigma',\tau'}\ar[r] &
\Gamma^{k+1}_{\sigma',\tau'}
}
\]
\oldpage[0\textsubscript{III}]{46}
para $(\sigma,\tau)\leq(\sigma',\tau')$. Entonces se verifica inmediatamente que
$\mathrm C^\bullet(A,B;\mathscr S^\bullet)
=\bigl(\mathrm C^h(A,B;\mathscr S^k)\bigr)_{(h,k)\in\mathbb Z\times\mathbb Z}$
es un \emph{bicomplejo} de grupos abelianos, y que
$\mathrm L^\bullet(A,B;\mathscr S^\bullet)
=\bigl(\mathrm L^h(A,B;\mathscr S^k)\bigr)$
es un subbicomplejo.
\end{env}

\begin{proposition}[11.8.9]
\label{0.11.8.9}
\label{0.11.8.9-fr}
La inyección canónica
$\mathrm L^\bullet(A,B;\mathscr S^\bullet)\to
\mathrm C^\bullet(A,B;\mathscr S^\bullet)$ define un isomorfismo para la cohomología de estos dos bicomplejos.
\end{proposition}

\begin{proof}
Pongamos $\mathrm C^{\bullet\bullet}=
\mathrm C^\bullet(A,B;\mathscr S^\bullet)$ y $\mathrm L^{\bullet\bullet}=
\mathrm L^\bullet(A,B;\mathscr S^\bullet)$ para simplificar, y observemos que, puesto que $\mathrm C^{hk}=\mathrm L^{hk}=0$ para $h<0$, las segundas
sucesiones espectrales de estos bicomplejos son \emph{regulares}
\sref{0.11.3.3}.
El homomorfismo
$\mathrm L^{\bullet\bullet}\to\mathrm C^{\bullet\bullet}$ da, por tanto,
un morfismo de sucesiones espectrales
${}''E(\mathrm L^{\bullet\bullet})\to
{}''E(\mathrm C^{\bullet\bullet})$
que, sobre los términos $E_2$, es
\[
\label{0.11.8.9.1-fr}
\tag{11.8.9.1}
H^p_{\mathrm{II}}(H^q_{\mathrm I}(\mathrm L^{\bullet\bullet}))
\longrightarrow
H^p_{\mathrm{II}}(H^q_{\mathrm I}(\mathrm C^{\bullet\bullet})).
\]
Sin embargo, para todo $k\in\mathbb Z$, \sref{0.11.8.3} muestra que
$H^q_{\mathrm I}(\mathrm L^{\bullet,k})\to
H^q_{\mathrm I}(\mathrm C^{\bullet,k})$
es \emph{biyectivo}.
La conclusión se deduce de \sref{0.11.1.5}.
\end{proof}

\begin{env}[11.8.10]
\label{0.11.8.10}
\label{0.11.8.10-fr}
Análogamente, con la notación de \sref{0.11.8.6}, existe un
homomorfismo canónico de bicomplejos
$\mathrm C^\bullet(A,B;\mathscr S^\bullet)\to
\mathrm P^\bullet(A,B;\mathscr S^\bullet)$
con la notación evidente.
El razonamiento de \sref{0.11.8.9}, utilizando ahora \sref{0.11.8.6}, muestra que
este homomorfismo induce asimismo un \emph{isomorfismo} en cohomología.
\end{env}

\subsection{Un lema sobre los complejos de tipo finito}
\label{subsection:0.11.9}
\label{subsection:0.11.9-fr}

\begin{proposition}[11.9.1]
\label{0.11.9.1}
\label{0.11.9.1-fr}
Sea $\boldsymbol{C}$ una categoría abeliana, y sean $\boldsymbol{K}'$ y $\boldsymbol{K}''$ subconjuntos
del conjunto de objetos de $\boldsymbol{C}$, con $\boldsymbol{K}''\subset\boldsymbol{K}'$, que satisfacen las
condiciones siguientes:
\begin{enumerate}
  \item[{\rm(i)}]
  Para todo $A'\in\boldsymbol{K}'$ y todo epimorfismo $u:A\to A'$ en
  $\boldsymbol{C}$, existen un objeto $B\in\boldsymbol{K}''$ y un morfismo $v:B\to A$ tales que $uv$ sea un epimorfismo.
\item[{\rm(ii)}] Para todo par de objetos $A\in\boldsymbol{K}'$, $B\in\boldsymbol{K}'$ y todo epimorfismo $u:A\to B$,
  el objeto $\operatorname{Ker}(u)$ pertenece a $\boldsymbol{K}'$.

  \item[{\rm(iii)}]
  El producto de dos objetos de $\boldsymbol{K}''$ pertenece a $\boldsymbol{K}''$.
\end{enumerate}

Sea $P_\bullet=(P_i)_{i\in\mathbb Z}$ un complejo en $\boldsymbol{C}$ tal
que $H_i(P_\bullet)\in\boldsymbol{K}'$ para todo $i$, y que exista $d$ para el cual $H_i(P_\bullet)=0$ cuando $i<d$.
Entonces existen un complejo $Q_\bullet=(Q_i)_{i\in\mathbb Z}$ en $\boldsymbol{C}$
tal que $Q_i\in\boldsymbol{K}''$ para todo $i$ y $Q_i=0$ para $i<d$, y un
homomorfismo de complejos
$u:Q_\bullet\to P_\bullet$
cuyo morfismo inducido
$H_\bullet(Q_\bullet)\to H_\bullet(P_\bullet)$
es un isomorfismo.
\end{proposition}

\begin{proof}
Demostremos primero la consecuencia siguiente de la condición~(i).

\medskip\noindent
(i \emph{bis})
\emph{Sea $u:C\to B$ un epimorfismo en $\boldsymbol{C}$, sea $A$ un
objeto de $\boldsymbol{K}'$, y sea $v:A\to B$ un morfismo en $\boldsymbol{C}$.
Entonces existen un objeto $D\in\boldsymbol{K}''$, un epimorfismo $u':D\to A$ y un
morfismo $v':D\to C$ tales que el diagrama}
\[
\xymatrix{
D\ar[r]^{u'}\ar[d]_{v'} & A\ar[d]^{v}\\
C\ar[r]_{u} & B
}
\tag{11.9.1.1}\label{0.11.9.1.1-fr}
\]
\emph{es conmutativo.}

Formemos el producto fibrado $C\times_B A$ en $\boldsymbol{C}$ y las proyecciones canónicas $p:C\times_B A\to C$, $q:C\times_B A\to A$, que hacen conmutativo el diagrama
\oldpage[0\textsubscript{III}]{47}
\[
\xymatrix{
C\times_B A\ar[r]^{q}\ar[d]_{\boldsymbol{p}} & A\ar[d]^{\boldsymbol{v}}\\
C\ar[r]_{\boldsymbol{u}} & B
}
\tag{11.9.1.2}\label{0.11.9.1.2-fr}
\]
es conmutativo.
El conúcleo de $q$ es el cociente de $A$ por $v^{-1}(u(C))$
([27, p.~1-12]).
Como $u$ es un epimorfismo, $u(C)=B$ y
$v^{-1}(u(C))=A$; por tanto, $q$ es un epimorfismo; basta entonces aplicar (i) al epimorfismo $q:C\times_B A\to A$, se obtienen un objeto
$D\in\boldsymbol{K}''$ y un morfismo $w:D\to C\times_B A$ tales que $qw$ sea un epimorfismo; se tomará $u'=qw$, $v'=pw$.

Dicho esto, para demostrar la proposición, procedamos por inducción. Supongamos, para un $i\geq d-1$, que para todo $j\leq i$, los objetos $Q_j$, los morfismos $d_j:Q_j\to Q_{j-1}$ y los morfismos $u_j:Q_j\to P_j$ de modo que $Q_j=0$ para $j<d$, que $d_{j-1}\circ d_j=0$ y $d_j\circ u_j=u_{j-1}\circ d_j$ para $j\leq i$; además, suponemos satisfechas las condiciones siguientes:
\begin{enumerate}
\item[$(\mathrm I_i)$]
  $Q_j\in\boldsymbol{K}''$ para $j\leq i$, y
  $B_j(Q_\bullet)\in\boldsymbol{K}'$ para $j<i$.

  \item[$(\mathrm{II}_i)$]
  Para $j<i$, el homomorfismo
  $H_j(Q_\bullet)\to H_j(P_\bullet)$
  inducido por $(u_k)_{k\leq i}$ es un \emph{isomorfismo}.

  \item[$(\mathrm{III}_i)$]
  El compuesto
  $v_i:Z_i(Q_\bullet)\to Z_i(P_\bullet)
\to H_i(P_\bullet)$
(donde la flecha de la izquierda es la restricción de $u_i$ y la flecha de la derecha el morfismo canónico) es un \emph{epimorfismo}.
\end{enumerate}

Observemos que, según (ii), $Z_i(Q_\bullet)$, núcleo del epimorfismo $Q_i\to B_{i-1}(Q_\bullet)$, pertenece a $\boldsymbol{K}'$ en virtud de la hipótesis $(\mathrm I_i)$. Se deduce además de (ii) que
$N_i=\operatorname{Ker}(v_i)$
pertenece a $\boldsymbol{K}'$, habida cuenta de la hipótesis $(\mathrm{III}_i)$. En virtud de (i \emph{bis}), existe un $Q'_{i+1}\in\boldsymbol{K}''$, un epimorfismo
$d'_{i+1}:Q'_{i+1}\to N_i$
y un morfismo $u'_{i+1}:Q'_{i+1}\to P_{i+1}$ tales que
\[
\xymatrix{
Q'_{i+1}\ar[r]^{d'_{i+1}}\ar[d]_{u'_{i+1}} &
N_i\ar[d]^{u_i}\\
P_{i+1}\ar[r]_{d_{i+1}} & B_i(P_\bullet)
}
\tag{11.9.1.3}\label{0.11.9.1.3-fr}
\]
es conmutativo.

El morfismo canónico
$Z_{i+1}(P_\bullet)\to H_{i+1}(P_\bullet)$
es un epimorfismo, y $H_{i+1}(P_\bullet)\in\boldsymbol{K}'$.
La condición~(i) proporciona, por tanto, un objeto $Q''_{i+1}\in\boldsymbol{K}''$ y un
morfismo
$u''_{i+1}:Q''_{i+1}\to Z_{i+1}(P_\bullet)$
tal que el compuesto
$Q''_{i+1}\to Z_{i+1}(P_\bullet)\to H_{i+1}(P_\bullet)$
es un epimorfismo.
Si $d''_{i+1}:Q''_{i+1}\to N_i$ igual a $0$, el diagrama
\[
\xymatrix{
Q''_{i+1}\ar[r]^{d''_{i+1}}\ar[d]_{u''_{i+1}} &
N_i\ar[d]^{u_i}\\
Z_{i+1}(P_\bullet)\ar[r]_{d_{i+1}} & P_i
}
\tag{11.9.1.4}\label{0.11.9.1.4-fr}
\]
\oldpage[0\textsubscript{III}]{48}
es conmutativo, siendo la flecha horizontal inferior $0$. Tomemos entonces $Q_{i+1}=Q'_{i+1}\times Q''_{i+1}$, que pertenece a $\boldsymbol{K}''$ en virtud de (iii), y
$d_{i+1}=d'_{i+1}+d''_{i+1}$, $u_{i+1}=u'_{i+1}+u''_{i+1}$.
Como $d_{i+1}(Q_{i+1})=d'_{i+1}(Q_{i+1})=N_i\subset Z_i(Q_\bullet)$, se tiene
$d_{i+1}\circ d_i=0$ y, con las notaciones usuales,
$B_i(Q_\bullet)=N_i$; esto demuestra
$(\mathrm I_{i+1})$.
La conmutatividad de \eqref{0.11.9.1.3} y
\eqref{0.11.9.1.4} da
$d_{i+1}\circ u_{i+1}=u_i\circ d_{i+1}$.
Por la definición de $N_i$, el morfismo
$H_i(Q_\bullet)=Z_i(Q_\bullet)/N_i\to H_i(P_\bullet)$
inducido por los $u_k$ ($k\leq i+1$) es el morfismo deducido de $v_i$ por paso a los cocientes, por tanto es un isomorfismo puesto que $v_i$ es un epimorfismo, de donde $(\mathrm{II}_{i+1})$. Finalmente, se tiene $Q''_{i+1}\subset Z_{i+1}(Q_\bullet)$ por definición; la elección de $u''_{i+1}$ muestra que
$v_{i+1}:Z_{i+1}(Q_\bullet)\to Z_{i+1}(P_\bullet)
\to H_{i+1}(P_\bullet)$
es un epimorfismo, porque su restricción a $Q''_{i+1}$ ya lo
es.
Así se cumple $(\mathrm{III}_{i+1})$, la inducción continúa y se deduce la
proposición.
\end{proof}

\begin{corollary}[11.9.2]
\label{0.11.9.2-fr}
Sean $A$ un anillo noetheriano (no necesariamente conmutativo), $P_\bullet=(P_i)_{i\in\mathbb Z}$ un complejo de $A$-módulos por la derecha. Se supone que los $H_i(P_\bullet)$ son $A$-módulos de tipo finito y que existe $d$ tal que $H_i(P_\bullet)=0$ para $i<d$. Entonces existe un complejo $Q_\bullet=(Q_i)_{i\in\mathbb Z}$ formado por $A$-módulos por la derecha libres de rango finito, tal que $Q_i=0$ para $i<d$, y un homomorfismo $u:Q_\bullet\to P_\bullet$ de complejos, tal que el homomorfismo $H_\bullet(Q_\bullet)\to H_\bullet(P_\bullet)$ correspondiente a $u$ sea biyectivo.
\end{corollary}

\begin{proof}
Se aplica (\hyperref[0.11.9.1-fr]{11.9.1}) tomando para $\boldsymbol{C}$ la categoría de los $A$-módulos por la derecha, para $\boldsymbol{K}'$ (resp. $\boldsymbol{K}''$) el conjunto de los $A$-módulos de tipo finito (resp. el conjunto de los $A$-módulos libres de rango finito); la verificación de las condiciones (i), (ii) y (iii) de (\hyperref[0.11.9.1-fr]{11.9.1}) es inmediata, habida cuenta de la hipótesis de que $A$ es noetheriano.
\end{proof}

\begin{env}[11.9.3]
\label{0.11.9.3}
\label{0.11.9.3-fr}
\emph{Observaciones}.
\begin{enumerate}
  \item[{\rm(i)}]
  Bajo las hipótesis de \sref{0.11.9.2}, supongamos además que los
  $P_i$ son $A$-módulos \emph{planos} a derecha.
  Entonces, para todo $A$-módulo a izquierda $M$, el homomorfismo de complejos
  $u\otimes 1:Q_\bullet\otimes_A M\to P_\bullet\otimes_A M$
  induce de nuevo un \emph{isomorfismo}
  $H_\bullet(Q_\bullet\otimes_A M)\xrightarrow{\sim}
H_\bullet(P_\bullet\otimes_A M)$ de la homología como lo veremos en el cap.~III.
\item[{\rm(ii)}] La conclusión de (\hyperref[0.11.9.2-fr]{11.9.2}) ya no es necesariamente exacta cuando no se supone $A$ \emph{noetheriano}; en efecto, al aplicarlo a un complejo reducido a $0$ salvo en un solo término, se concluiría que todo $A$-módulo por la izquierda finitamente engendrado admite una \emph{resolución} por
  módulos libres finitamente engendrados, lo que es falso en general
(cf. Bourbaki, \emph{Alg. comm.}, cap.~I, \S~2, ejerc.~6).

  Sin embargo, es posible reemplazar la hipótesis de que $A$ es
  noetheriano por la suposición de que los $H_i(P_\bullet)$ admiten
  $\infty$-presentaciones finitas (cf.\ Capítulo~IV).
\end{enumerate}
\end{env}

\subsection{Característica de Euler--Poincar\'e de un complejo de módulos de longitud finita}
\label{subsection:0.11.10}
\label{subsection:0.11.10-fr}

\begin{env}[11.10.1]
\label{0.11.10.1}
\label{0.11.10.1-fr}
Sea $A$ un anillo, no necesariamente conmutativo, y sea
\[
\label{0.11.10.1.1-fr}
\tag{11.10.1.1}
M^\bullet:0\longrightarrow M^0\longrightarrow M^1
\longrightarrow\ldots\longrightarrow M^n\longrightarrow0
\]
un \emph{complejo} de $A$-módulos a izquierda de \emph{longitud finita}.
La \emph{característica de Euler--Poincar\'e} de este complejo es el
entero
\oldpage[0\textsubscript{III}]{49}
\[
\label{0.11.10.1.2-fr}
\tag{11.10.1.2}
\chi(M^\bullet)=\sum_{i=0}^{n}(-1)^i\operatorname{long}M^i.
\]
\end{env}

\begin{proposition}[11.10.2]
\label{0.11.10.2}
\label{0.11.10.2-fr}
Para todo complejo finito $M^\bullet$ de $A$-módulos a izquierda de longitud
finita,
$\chi(M^\bullet)=\chi(H^\bullet(M^\bullet))$
donde $H^\bullet(M^\bullet)$ se considera como un complejo con diferencial
nulo.
En particular, si la sucesión \eqref{0.11.10.1.1} es exacta, entonces
$\chi(M^\bullet)=0$.
\end{proposition}

\begin{proof}
Pongamos para abreviar $B^i=B^i(M^\bullet)$, $Z^i=Z^i(M^\bullet)$, $H^i=H^i(M^\bullet)=Z^i/B^i$; los $B^i$, $Z^i$ y $H^i$ tienen todos
longitud finita.
Las sucesiones exactas
\[
0\longrightarrow B^i\longrightarrow Z^i\longrightarrow H^i
\longrightarrow0
\]
y
\[
0\longrightarrow Z^i\longrightarrow M^i\longrightarrow B^{i+1}
\longrightarrow0
\]
dan
\[
\operatorname{long}(Z^i)=\operatorname{long}(H^i)+
\operatorname{long}(B^i)
\]
y
\[
\operatorname{long}(M^i)=\operatorname{long}(Z^i)+
\operatorname{long}(B^{i+1})
\]
Por consiguiente,
\[
\operatorname{long}(M^i)-\operatorname{long}(H^i)=
\operatorname{long}(B^{i+1})+\operatorname{long}(B^i)
\]
Multipliquemos esta igualdad por $(-1)^i$ y hagamos la suma de las relaciones obtenidas para $0\leq i\leq n$.
Como $B^0=B^{n+1}=0$, se deduce la igualdad deseada.
\end{proof}

\begin{corollary}[11.10.3]
\label{0.11.10.3}
\label{0.11.10.3-fr}
Sea $E=(E_r^{pq})$ una sucesión espectral en la
categoría de los módulos sobre un anillo $A$.
Se supone que los $E_2^{pq}$ son $A$-módulos de longitud finita y que solo hay un número finito de pares $(p,q)$ tales que $E_2^{pq}\neq0$. Entonces las características de Euler-Poincaré
$\chi(E_r^{(\bullet)})$
de todos los complejos
$E_r^{(\bullet)}=(E_r^{(n)})_{n\in\mathbb Z}$
de \sref{0.11.1.1} son iguales.
Si, además, $E$ es débilmente convergente y si se pone $E_\infty^{(n)}=\bigoplus_{p+q=n}E_\infty^{pq}$ para todo $n\in\mathbb Z$, se tiene también $\chi(E_\infty^{(\bullet)})=\chi(E_2^{(\bullet)})$,
donde
$E_\infty^{(\bullet)}=(E_\infty^{(n)})_{n\in\mathbb Z}$
se considera como un complejo con diferencial nulo.
\end{corollary}

\begin{proof}
Observemos primero que si $E_2^{pq}=0$, entonces
$E_r^{pq}=0$ para $2\leq r\leq+\infty$.
Así, todo complejo $E_r^{(\bullet)}$ es finito y está formado por
$A$-módulos de longitud finita; la relación $\chi(E_r^{(\bullet)})=\chi(E_{r+1}^{(\bullet)})$ para todo $r$ finito resulta, pues, de (\hyperref[0.11.10.2-fr]{11.10.2}) y de la isomorfía entre $H^\bullet(E_r^{(\bullet)})$ y $E_{r+1}^{(\bullet)}$ (como complejos con derivación trivial). La hipótesis de que los $E_r^{pq}$ tienen longitud finita, para cada par
$(p,q)$, las sucesiones
$(B_k(E_2^{pq}))_{k\geq2}$ y $(Z_k(E_2^{pq})_{k\geq2}$ son estacionarias;
la hipótesis de que $E$ es débilmente convergente y de que $E_2^{pq}=0$ salvo para un número finito de pares $(p,q)$, implica, por tanto,
que algún entero $r\geq2$ satisface
$E_\infty^{pq}=E_r^{pq}$
para todo par $(p,q)$.
Se deduce la aserción relativa a
$\chi(E_\infty^{(\bullet)})$.
\end{proof}
